% !TEX program = pdflatex
% Elsevier elsarticle format
% Suitable for: Robotics and Autonomous Systems / Computers & Electrical Engineering
\documentclass[preprint,12pt,3p]{elsarticle}

\usepackage{amsmath,amssymb,amsthm,bm}
\usepackage{graphicx}
\usepackage{booktabs}
\usepackage{array}
\usepackage{multirow}
\usepackage{algpseudocode}
\usepackage{algorithm}
\usepackage{url}
\usepackage[hidelinks,plainpages=false]{hyperref}
\hypersetup{pageanchor=false}
\usepackage{lineno}
\IfFileExists{microtype.sty}{\usepackage[final]{microtype}}{}
\emergencystretch=2em

\graphicspath{{figs/}}
\newtheorem{theorem}{Theorem}
\newtheorem{proposition}{Proposition}
\newtheorem{corollary}[theorem]{Corollary}
\newtheorem{lemma}[theorem]{Lemma}
\theoremstyle{definition}
\newtheorem{definition}{Definition}
\newtheorem{remark}[theorem]{Remark}
\newenvironment{bproof}{\begin{proof}}{\end{proof}}

\newcommand{\conv}{\operatorname{conv}}
\newcommand{\diam}{\operatorname{diam}}
\newcommand{\MEC}{\operatorname{MEC}}
\newcommand{\Rmin}{R_{\min}}
\newcommand{\Rmax}{R_{\max}}
\newcommand{\bA}{\mathcal{A}}
\newcommand{\bL}{\mathcal{L}}
\newcommand{\bP}{\mathcal{P}}
\newcommand{\bF}{\mathcal{F}}
\newcommand{\bE}{E}
\newcommand{\R}{\mathbb{R}}

\journal{Robotics and Autonomous Systems}

\begin{document}

\begin{frontmatter}

\title{Certifiable Search and Neutralization of Radio Emitters with Unknown
Directionality Using a Bearing-Only Mobile Sensor}

\author[a]{Changxuan Cao\corref{cor1}}
\ead{2023211904@ecut.edu.cn}
\author[a]{Zexu Ouyang}
\ead{2023211863@ecut.edu.cn}
\cortext[cor1]{Corresponding author.}
\address[a]{East China University of Technology, Nanchang, Jiangxi 330013, China}

\begin{abstract}
A mobile platform must find and neutralize an unknown number of radio-frequency emitters in a
bounded region using only noisy bearings, and must certify that none was missed. The emitter count
and reception radii are unknown, and an emitter may radiate into only a half-plane of unknown
orientation, so a silent channel carries no distance information. We derive a necessary and
sufficient condition for a network to detect an emitter of arbitrary orientation,
$G\in\conv(S_G)$, which forces the network beyond the region. We build a certifiable search
framework in which every channel ends cleared or proved empty, making termination decidable rather
than a timeout, and cast network design as a covering-salesman problem with a rigorous point-count
lower bound; a two-layer construction needs $25$ points on a route $42\%$ shorter than a
single-layer lattice. On $200$ Monte-Carlo episodes it clears $100\%$ of emitters at
$359.7\pm44.6$~s each; in a separate $30$-seed comparison the best of four baselines clears $51\%$
with $12.9$ times the budget and none certifies completion in the directional scenario. Ablating
the certificate leaves what is cleared unchanged but cuts mission time $4.8$-fold. Under multipath
a $q$-relaxed evidence set removes the false absence certificates the hard outer approximation
produces; sound and complete certification under $q$ outliers provably needs a detection redundancy
$r>q$ that covering networks lack.
\end{abstract}


\begin{keyword}
Bearing-only localization \sep Guaranteed search \sep Covering salesman problem \sep Directional emitters \sep Radio-frequency interference \sep Search and neutralize \sep Mobile sensors
\end{keyword}

\end{frontmatter}

\section{Introduction}

Radio-frequency (RF) interference is a first-order operational problem in
spectrum management, aviation safety and wireless network operation. Once a monitoring
infrastructure has narrowed an interferer down to a region of a few kilometres, the final
localization and the physical removal of the source are still performed by a human operator who
walks through the area with a handheld direction finder, interprets the evolution of field
strength and bearing, and finally locates and switches off the device. Automating this last,
dangerous and time-consuming step with a mobile robot is the problem considered here.

The task is deceptively simple to state: a mobile platform enters a bounded circular region and
must neutralize every emitter inside it as quickly as possible. The platform carries a direction
finder with a known angular uncertainty and a short-range optical/laser tool that can neutralize
an emitter once the platform is within a few tens of metres of it. What makes the problem
genuinely hard is the combination of four uncertainties.

\begin{itemize}
\item \textbf{The number of emitters is unknown.} Neither the count nor the set of active
channels is revealed. Consequently a search cannot stop when a counter reaches a target value;
it must stop when every possible emitter has either been neutralized or \emph{proved}
non-existent.
\item \textbf{No distance information.} The sensor measures bearings only. A single bearing
constrains the emitter to a wedge but says nothing about its range, and the effective reception
radius differs from emitter to emitter and is not reported.
\item \textbf{Unknown directionality.} An emitter may be omni-directional or may radiate only
into a half-plane whose orientation is unknown. In the latter case a non-detection is ambiguous:
it may mean that no emitter exists, that the platform is beyond the reception radius, or that the
platform lies outside the illuminated half-plane.
\item \textbf{Absence of a distance measurement in the certificate.} Because of the previous
point, the standard ``coverage'' argument --- if every point of the region is within reception
range of some measurement location, no emitter can hide --- is no longer sufficient: an emitter
may be well within range and still be invisible from a given location.
\end{itemize}

\subsection{Contributions}

This paper makes the following contributions.

\begin{enumerate}
\item \textbf{A detectability condition for orientation-unknown emitters}
(Theorem~\ref{thm:detect}). We show that a set of measurement locations detects an emitter of
\emph{arbitrary} orientation at position $G$ if and only if $G$ lies in the convex hull of the
measurement points that are within the (worst-case) reception radius of $G$. The condition is
the dual of classical coverage: coverage asks for \emph{some} point nearby, whereas detection of
an arbitrarily oriented emitter asks for points \emph{surrounding} it. A direct corollary is that
the measurement network must extend beyond the region of interest, since an emitter on the
boundary radiating outwards is invisible from every point inside.
\item \textbf{A certifiable search framework} (Section~\ref{sec:method}). Every channel is driven
to one of two terminal states: \emph{cleared} or \emph{certified absent}. Certified absence is
obtained from the detectability condition, so the mission terminates with a proof rather than a
timeout. Ablating the certificate leaves the number of neutralized emitters unchanged in both
scenarios, but in S2 it increases mission time by a factor of $4.8$, because there the census
certificate is the only available route to an absence proof (in S1 the raster path already resolves
empty channels, so removing the certificate changes nothing).
\item \textbf{A network design with guarantees} (Section~\ref{sec:network}). We cast the network
design as a covering-salesman problem, prove a six-point lower bound for the omni-directional case
and exhibit a seven-point layout that satisfies the coverage requirement, and derive a route-length
lower bound. For the directional case we show that a two-layer construction --- a dense inner
lattice plus a ring placed just outside the region --- satisfies the detectability condition --- analytically almost everywhere and numerically on the residual --- with $25$ points and a $17.7$~km route, $42\%$ shorter than the natural single-layer lattice.
\item \textbf{A robustness analysis for non-ideal channels} (Section~\ref{sec:robust}). We show
that the hard outer approximation is \emph{unsafe} once spurious bearings occur --- a $2\%$
spurious-bearing rate produces false absence certificates --- and that a $q$-relaxed intersection
with a separated absence path eliminates them. We further prove that sound \emph{and} complete
absence certification under $q$ outliers requires a detection redundancy $r>q$ that the covering
networks used here do not possess, which we state as an open network-design problem.
\item \textbf{An experimental separation of inference from planning}
(Section~\ref{sec:planning}). Holding the evidence sets and the certificate fixed and changing one
planning decision at a time, we show that the inference layer decides \emph{whether} the mission
succeeds (up to $-58\%$ when removed) while the planning layer decides \emph{how fast} (up to
$+78\%$ when degraded), and that the two axes are independent.
\item \textbf{A sample-level array-processing surrogate for hardware}
(Section~\ref{sec:rf}). We replace the abstract bearing model by a complete direction-finding
chain --- specular multipath, uniform circular array, snapshot covariance, MUSIC --- and use it to
test the $\delta=1^\circ$ assumption rather than postulate it. The assumption holds in a
line-of-sight channel ($99$th percentile error $0.85^\circ$ at $12$~dB) and fails under multipath;
we further show that antenna \emph{aperture}, not integration time, is what restores it, and that
with an adequate aperture and a calibrated $\delta$ the full method clears $100\%$ of emitters with
zero false certificates.
\item \textbf{An extension to obstructed workspaces} (Section~\ref{sec:obstacles}). The
detectability condition becomes an enclosure condition over line-of-sight-reachable network points;
a disk-optimized network fails the condition at $41\%$ of the free-space positions of an
obstructed workspace, and greedy shadow-filling restores it to $97.7\%$ at roughly twice the
number of points.
\item \textbf{A thorough empirical study} (Section~\ref{sec:experiments}). We compare against four
baselines that enjoy exactly the same information, report paired sign tests with bootstrap
confidence intervals, ablate six components, sweep five generalization axes, and compare against
an omniscient offline planner as a competitiveness benchmark.
\end{enumerate}

%=====================================================================
\section{Related Work}

\subsection{Bearing-only and emitter localization}
Bearing estimation itself is a mature subject: the MUSIC subspace method of
Schmidt~\cite{schmidt1986} and the array-processing review of Krim and
Viberg~\cite{krim1996} are the standard references for how the angle is obtained, which we take as
given. Bearing-only (angle-of-arrival) localization has a long history in passive surveillance;
Torrieri~\cite{torrieri1984} established the statistical theory of passive location systems and
Gavish and Weiss~\cite{gavish1992} analysed the accuracy of bearing-only target location, showing
that the achievable covariance is governed by the sensor--target geometry. Bishop
\emph{et al.}~\cite{bishop2010} derived optimality conditions for sensor--target localization
geometries, and Mart\'inez and Bullo~\cite{martinez2006} optimized sensor placement and motion
coordination for target tracking. These results characterize the \emph{accuracy} of an estimate
for a \emph{known number} of targets and, crucially, they assume that the target is detected
whenever it is within range. Our setting differs in that detection itself is uncertain because
the radiation pattern is unknown, and in that the number of emitters is unknown, so the decision
problem is not only ``where to measure'' but also ``when may I stop''.

Recent work has advanced both halves of this problem on real platforms. Peng \emph{et al.}
\cite{peng2024} shape the platform trajectory so as to maximise the observability of a
bearing-only estimate while calibrating sensor bias; Kudyba \emph{et al.}~\cite{kudyba2024}
build a radio map from an autonomous aerial drone by Bayesian optimization; and Kim
\emph{et al.}~\cite{kim2025} localize a radio source with a Gaussian-mixture filter on a
mobile robot. These systems estimate \emph{where} an emitter is, and how accurately, and they
are demonstrably effective at that. What none of them provides is a statement that \emph{no}
emitter remains, which is the question this paper is organized around.

When the waveform is accessible, time-difference and frequency-difference measurements provide
range information that bearing-only sensing lacks; the classical closed-form estimators of
Chan and Ho~\cite{chan1994} and Ho and Xu~\cite{ho2004} are the standard tools, and
aerial platforms are the usual carriers for such arrays. All of this work assumes that the
emitter is received by every observing platform; none of it addresses an emitter that radiates
into a half-plane of unknown orientation, which is the difficulty we concentrate on.

\subsection{Set-membership estimation and guaranteed inference}
\label{sec:relset}
Our evidence set $E_c$ is a \emph{set-membership} (also called bounded-error or guaranteed)
state estimator: rather than maintaining a probability density over the emitter position, it
maintains the set of positions consistent with every measurement and with a hard bound on the
measurement error. The theory of such estimators is well developed. Milanese and
Vicino~\cite{milanese1991} gave the foundational overview for dynamic systems with set-membership
uncertainty, and Jaulin \emph{et al.}~\cite{jaulin2001} developed interval analysis as the
computational substrate. Because exact sets can be expensive to propagate, a large body of work
approximates them by tractable families: parallelotopes~\cite{chisci1996}, ellipsoids and
zonotopes~\cite{alamo2005,combastel2003}, and constrained zonotopes~\cite{scott2016}. The
bounded-error view has recently been carried to multi-vehicle target localization under
partial visibility: Ibenthal \emph{et al.}~\cite{ibenthal2023} localize partially hidden
moving targets with a fleet of UAVs using bounded-error estimation. That work asks where a
\emph{known} set of targets is; here the open question is whether any emitter remains, and it
is the absence question --- not the estimation accuracy --- that drives the network design and
the termination rule.

Our setting differs from this literature in two respects that turn out to be essential. First,
the update here is \emph{not} a linear measurement with bounded additive noise but a disjunctive
sensor model \eqref{eq:sensor}: a reading either yields an angular constraint, or a
\texttt{near} constraint, or --- for a directional emitter --- no constraint at all. The
corresponding sets are convex cones intersected with disks and their complements, which are
non-convex even though each individual update is convex; we therefore use a raster outer
approximation rather than a zonotope. Second, and more importantly, the existing literature is
concerned with the \emph{accuracy} of the estimate, whereas we are concerned with what the
estimate certifies. A bounded set that has become empty is a proof of absence, and the whole
mission-termination argument of Section~\ref{sec:theory} rests on that observation rather than on
the size of the set.

A second, complementary line is the placement of the measurement locations themselves so that the
estimator is well conditioned. We address this for the second measurement in
Theorem~\ref{thm:second}, but a related body of work optimizes the \emph{error covariance} rather
than a feasibility set~\cite{bishop2010,martinez2006}. The distinction matters when the error is
bounded rather than Gaussian: a geometry that is optimal for the covariance need not be the one
that shrinks the feasible set fastest, and the latter is what determines when the clearance
criterion of Theorem~\ref{thm:clear} can be applied.

\subsection{Coverage path planning}
Coverage path planning is surveyed by Choset~\cite{choset2001} and by Galceran and
Carreras~\cite{galceran2013}; recent reviews cover heuristic and classical
algorithms~\cite{tan2021} and unmanned-aerial-vehicle applications~\cite{cabreira2019}. The
boustrophedon (lawnmower) decomposition~\cite{choset1998} is the canonical complete coverage
strategy and is the natural baseline for our problem; the general principles are collected
in~\cite{choset2005}. Classical coverage guarantees that every point of the workspace is visited
within the footprint of a sensor whose detection model is \emph{isotropic}: being close is
sufficient. Our detectability condition shows that when the source may radiate into an unknown
half-plane, ``being close'' is not sufficient, and the network must instead \emph{enclose} every
point. To our knowledge this enclosure condition has not been stated for emitter search.

For multi-robot and target-oriented variants, Robin and Lacroix~\cite{robin2016} and Khan
\emph{et al.}~\cite{khan2018} survey target detection and tracking, and cooperative air--ground
surveillance was demonstrated by Grocholsky \emph{et al.}~\cite{grocholsky2006}. These works
assume probabilistic detection models that do not provide the deterministic completeness
certificate required here; the certificate we derive is a property of the geometry of the
network, not of a likelihood.

\subsection{Covering salesman and network design}
The problem of choosing a shortest tour such that the union of neighbourhoods of the visited
points covers a region is the covering salesman problem; the generalized variant was formalized by
Golden \emph{et al.}~\cite{golden2012}. We use this formulation, add a rigorous lower bound on the
number of points from a boundary-arc argument, and show that within the design family we consider --- a lattice plus an outer ring --- the best construction is a two-layer network rather than a single extended lattice. Bounds for covering a disk by equal circles are
classical~\cite{nica2011,jung1901,toth1998}, and the clustering view of the same problem is the
minimum-enclosing-circle and $k$-centre family~\cite{gonzalez1985,preparata1985}. The tour itself
is a Euclidean travelling-salesman problem, for which the asymptotic constant we use in
Section~\ref{sec:experiments} goes back to Beardwood, Halton and
Hammersley~\cite{beardwood1959}; modern treatments are given by
Applegate \emph{et al.}~\cite{applegate2006}.

\subsection{Source seeking and search theory}
When a scalar field is measurable, gradient climbing and Bayesian search strategies such as
infotaxis~\cite{vergassola2007} are highly effective, and stochastic source-seeking laws with
provable convergence exist when the field is smooth~\cite{azuma2012}. Search theory provides
optimal allocation of search effort for a known prior~\cite{stone2007,koopman1980}, and the
probabilistic-robotics view of the same problem is given by
Thrun \emph{et al.}~\cite{thrun2005}. The recent literature continues in this direction,
replacing hand-designed heuristics by learned policies for adaptive informative path
planning~\cite{popovic2024}; multi-robot variants now scale to teams through deep reinforcement
learning~\cite{vashisth2026}. The objective remains expected information gain under a prior, so
the same observation applies. These formulations optimize an expected reward under a
prior; they do not provide the deterministic completeness certificate required in our
formulation. Our problem differs in two ways:
the sensor does not return any field strength or range, only a bearing, and the requirement is a
\emph{certificate} of completeness rather than a probabilistic statement. The related literature on guaranteed search is surveyed by Chung, Hollinger and
Isler~\cite{chung2011}; the guarantees obtained there rest on omnidirectional sensing, in contrast
to the half-plane emitters considered here.

\subsection{Directional sensor networks}
Coverage in directional sensor networks, where each sensor has a bounded field of view with an
orientation, has been surveyed by Guvensan and Yilmaz~\cite{guvensan2011}. The typical question
there is how to orient a \emph{given} set of sensors to maximize coverage. Our question is the
inverse: given that the \emph{source} orientation is unknown and adversarial, where must the
mobile sensor measure so that every position is detectable for every orientation?

%=====================================================================
\section{Problem Formulation}
\label{sec:problem}

\subsection{Workspace and emitters}
The workspace $\bA=\{p\in\R^2:\|p\|\le R\}$ is a disk of radius $R=1800$~m centred at the
origin. An unknown number $k$ of emitters is present, $k\in[k_{\min},k_{\max}]=[10,16]$. Emitter
$i$ occupies a distinct \emph{channel} $c_i\in\{1,\dots,N_c\}$, $N_c=20$; channels do not change
during a mission. Emitter $i$ is characterized by
\begin{itemize}
\item its position $G_i\in\bA$;
\item its \emph{effective reception radius} $\rho_i\in[\Rmin,\Rmax]=[1000,1500]$~m, unknown and
different for each emitter;
\item its \emph{radiation pattern}: either omni-directional, or directional with an unknown
orientation $d_i\in[0,2\pi)$ illuminating the half-plane
$H(G_i,d_i)=\{p:(p-G_i)\cdot u(d_i)\ge0\}$, where $u(\alpha)=(\cos\alpha,\sin\alpha)$.
\end{itemize}

\subsection{Sensor model}
\label{sec:sensor}
When the platform measures channel $c$ at position $p$, it obtains
\begin{equation}
z=\begin{cases}
\texttt{near}, & \|p-G_c\|\le r_{\mathrm{near}}\ \text{and}\ p\in H(G_c,d_c),\\
\theta_c(p)+\varepsilon, & \|p-G_c\|\le\rho_c\ \text{and}\ p\in H(G_c,d_c),\\
\texttt{no\_signal}, & \text{otherwise},
\end{cases}
\label{eq:sensor}
\end{equation}
where $r_{\mathrm{near}}=5$~m is a saturation radius inside which no bearing can be resolved
(the platform may neutralize directly), and
$\theta_c(p)=\operatorname{atan2}(G_{c,y}-p_y,\;G_{c,x}-p_x)$ is the true bearing. The measurement
noise $\varepsilon$ is bounded, $|\varepsilon|\le\delta$, with $\delta=1^\circ$ by default. A key
structural assumption, inherited from the physical setting, is that \emph{the error at a given
location is deterministic}: repeating a measurement at the same point returns the same reading.
Only moving to a different location produces an independent error. Consequently averaging
repeated measurements at one location yields no improvement, and any localization procedure must
exploit \emph{spatial} diversity.

Note that \eqref{eq:sensor} deliberately does \emph{not} provide range. A reading of
\texttt{no\_signal} is triply ambiguous: it can mean that channel $c$ is empty, that the platform
is beyond $\rho_c$, or that the platform is outside the illuminated half-plane.

\subsection{Motion and neutralization}
The platform moves in straight lines at speed $v=5$~m/s. A measurement costs $t_m=5$~s plus
$t_s=1$~s if the channel is changed. Neutralization is attempted by a short-range optical/laser
tool: it costs $t_o=3$~s to attempt optical localization and, if an emitter of the requested
channel is within $r_{\mathrm{clr}}=20$~m, an additional $t_c=2$~s to neutralize it. Crucially the
neutralization succeeds \emph{irrespective of the radiation pattern}, so directionality affects
only discovery, never removal.

\subsection{Objective}
Let $T$ denote the total mission time, i.e. the sum of travel, channel switching, measurement and
neutralization times, and let $n$ be the number of neutralized emitters. The objective is
\begin{equation}
\min \; T \qquad\text{subject to}\qquad n=k\ \ \text{for every realization,}
\label{eq:objective}
\end{equation}
together with the requirement that the platform terminate by its own decision. Because $k$ is
unknown, the second constraint cannot be verified by counting; it must be established by
exhausting the space of possibilities. This is the sense in which our method is \emph{certifiable}:
a mission ends only when every channel is either neutralized or proved empty.

\subsection{Assumptions and scope}
(A1) No obstacles; the workspace is flat and the platform and emitters are coplanar.
(A2) Errors at different locations are independent; errors at the same location are identical
(see Section~\ref{sec:sensor}).
(A3) A channel hosts at most one emitter at a time. Multiple collocated emitters sharing one
channel are excluded from the model; Remark~\ref{rem:cochannel} states what changes if that
assumption is dropped.
(A4) The platform may also measure outside $\bA$; this is permitted by the application and is
exploited in Section~\ref{sec:network}.
(A5) A lower bound $\Rmin$ on the effective reception radius, and a bound $\delta$ on the bearing
error, are known.
(A6) In the baseline model every reading is truthful; Section~\ref{sec:robust} relaxes this to
``at most $q$ spurious readings per channel'' and states what that costs.

\noindent These assumptions delimit the claim we make. In particular, ``unknown
directionality'' means \emph{omni-directional, or an ideal half-plane of unknown orientation}; it does not
cover arbitrary lobes, polarization effects, or a source that switches itself off intermittently.
The absence certificate of Theorem~\ref{thm:cert} is a statement about this model, and its
validity depends on (A5) and (A6) holding.

%=====================================================================
\section{Theory}
\label{sec:theory}

\subsection{Evidence sets and conservative outer approximation}

For each channel $c$ we maintain an \emph{evidence set} $E_c\subseteq\bA$, the set of positions
that are still consistent with every reading obtained so far for that channel. The true emitter,
if it exists, always belongs to $E_c$; this invariant is what makes the method certifiable.

\begin{lemma}[Wedge constraints]
\label{lem:wedge}
Let $f(X;\alpha)=(X-S)\times u(\alpha)=(x-S_x)\sin\alpha-(y-S_y)\cos\alpha$. Then
\begin{equation}
W(S,\theta)=\{X:f(X;\theta-\delta)\le0\}\ \cap\ \{X:f(X;\theta+\delta)\ge0\}
\label{eq:wedge}
\end{equation}
is exactly the closed convex cone with apex $S$, central bearing $\theta$ and half-angle
$\delta$, i.e.\ $W(S,\theta)=\{S+t\,u(\alpha)\}$ with $t\ge0$ and
$\alpha\in[\theta-\delta,\theta+\delta]$.
\end{lemma}
\begin{proof}
For $X=S+t\,u(\beta)$ with $t>0$ we have $f(X;\alpha)=t\sin(\alpha-\beta)$. Hence
$f(X;\theta-\delta)\le0\iff\sin(\theta-\delta-\beta)\le0\iff\beta\ge\theta-\delta$ and
$f(X;\theta+\delta)\ge0\iff\beta\le\theta+\delta$ for $\beta$ in the relevant branch; the
negative ray ($t<0$, i.e. $\beta=\theta+\pi$) violates the first inequality. Conversely every
point satisfying both inequalities lies on a ray $S+t\,u(\beta)$, $t\ge0$, with
$\beta\in[\theta-\delta,\theta+\delta]$.
\end{proof}

When a bearing reading is received, we intersect $E_c$ with $W(p,\bar\theta)$ and with the disk
$\bar B(p,\Rmax)$; a \texttt{near} reading yields $E_c\cap\bar B(p,r_{\mathrm{near}})$; an
omni-directional \texttt{no\_signal} reading removes the open disk $B(p,\Rmin)$, because an
omni-directional emitter would certainly have been heard inside that radius. For a directional
emitter \texttt{no\_signal} admits no distance inference and leaves $E_c$ unchanged. A failed
neutralization removes $B(p,r_{\mathrm{clr}})$, which is valid for both pattern types.

Because $E_c$ is represented on a $10$~m raster, all set operations are implemented as
\emph{conservative outer approximations}: a grid cell is discarded only when the whole cell is
provably inconsistent. The underlying primitives --- convex-hull computation, half-plane clipping
and point-in-polygon tests --- are standard~\cite{preparata1985}. Denote by $\eta$ the
circumradius of a cell.

\begin{lemma}[Safety of the outer approximation]
\label{lem:outer}
If $G_c$ lies in cell $\kappa$, then $\kappa$ is never removed by any of the update rules of
Table~\ref{tab:ops}.
\end{lemma}
\begin{proof}
For each rule the cell is removed only if every point of the cell satisfies the negated
constraint with a margin of at least $\eta$. Since $\|X-x_\kappa\|\le\eta$ for all $X$ in cell
$\kappa$ (centre $x_\kappa$), the true position $G_c$ satisfies the constraint strictly whenever
the cell is removed, contradicting $G_c\in\kappa$.
\end{proof}

\begin{table}[t]
\caption{Evidence update rules ($p$ is the measurement position).}
\label{tab:ops}
\centering
\small
\begin{tabular}{lll}
\toprule
Reading & Omni-directional & Directional \\
\midrule
direction ($\bar\theta$) & $E\cap W(p,\bar\theta)\cap\bar B(p,\Rmax)$ & same \\
\texttt{near} & $E\cap\bar B(p,r_{\mathrm{near}})$ & same \\
\texttt{no\_signal} & $E\setminus B(p,\Rmin)$ & $E$ (unchanged) \\
failed neutralization & $E\setminus B(p,r_{\mathrm{clr}})$ & same \\
successful neutralization & channel closed & channel closed \\
\bottomrule
\end{tabular}
\end{table}

\begin{theorem}[Clearance criterion]
\label{thm:clear}
Let $c$ be an unresolved channel with evidence set $E_c\ne\varnothing$, let $x^\ast$ be the
centre of the minimum enclosing circle (MEC) of $E_c$, computed in linear time by the randomized
algorithm of Welzl~\cite{welzl1991} (see also~\cite{megiddo1983}), and let
$r^\ast=\max_{X\in E_c}\|X-x^\ast\|$. If $r^\ast\le r_{\mathrm{clr}}-\eta$ then a single
neutralization attempt at $x^\ast$ succeeds with certainty.
\end{theorem}
\begin{proof}
By Lemma~\ref{lem:outer} the true emitter $G_c$ lies in some non-removed cell $\kappa$, hence
$\|G_c-x_\kappa\|\le\eta$ and $x_\kappa\in E_c$. Therefore
$\|G_c-x^\ast\|\le\|G_c-x_\kappa\|+\|x_\kappa-x^\ast\|\le\eta+r^\ast\le r_{\mathrm{clr}}$, and
neutralization succeeds because it depends only on distance, not on the radiation pattern.
\end{proof}

Theorem~\ref{thm:clear} is the operational stop criterion: it converts a hard geometric question
(``where exactly is the emitter?'') into a decidable one (``is the MEC radius below a threshold?'').
In the implementation we additionally use an exact polygonal representation obtained by
intersecting the bearing wedges, which is tighter than the raster and has no $\eta$ penalty.

\subsection{Placing the second measurement}

Consider a single emitter about which one bearing $\theta_1$ has been obtained at $S_1$. The
consistent set is
$\bF_1=\bA\cap W(S_1,\theta_1)\cap\{r_{\min}\le\|X-S_1\|\le\Rmax\}$, where $r_{\min}=r_{\mathrm{near}}$
excludes the saturated zone. A second measurement at $S_2$ is only useful if the emitter is
\emph{certain} to be heard from $S_2$; otherwise the second reading may be an uninformative
\texttt{no\_signal}. Writing $S_2=S_1+d\,u(\theta_1+\psi)$, the decisive quantity is
$\max_{X\in\bF_1}\|S_2-X\|$, and it admits a closed form.

\begin{theorem}[Feasible set of the second measurement; omni-directional emitter]
\label{thm:second}
Assume that the emitter is omni-directional, so that reception depends on distance alone; the
directional case is discussed in Remark~\ref{rem:dirsecond}.
With $X=S_1+r\,u(\theta_1+\beta)$, $\beta\in[-\delta,\delta]$, and $\Theta=|\psi|+\delta$,
\begin{equation}
\max_{X\in\bF_1}\|S_2-X\|=\sqrt{\max\{d^2+r_{\min}^2-2dr_{\min}\cos\Theta,\;
d^2+\Rmax^2-2d\Rmax\cos\Theta\}}.
\label{eq:maxdist}
\end{equation}
Imposing this quantity $\le \Rmin$ is equivalent to the two quadratic inequalities
\begin{align}
d^2-2\Rmax\cos\Theta\,d+(\Rmax^2-\Rmin^2)&\le0,\label{eq:c1}\\
d^2-2r_{\min}\cos\Theta\,d+(r_{\min}^2-\Rmin^2)&\le0.\label{eq:c2}
\end{align}
Consequently the admissible second measurement positions form a curved rectangle
\begin{equation}
\bP_2=\{S_1+d\,u(\theta_1+\psi):|\psi|\le\psi_{\max},\ d\in[d_1(\psi),d_2(\psi)]\}
\end{equation}
where $d_{1,2}$ are the roots of \eqref{eq:c1} intersected with the upper root of \eqref{eq:c2},
and $\psi_{\max}$ is the largest $\psi$ for which those two intervals still intersect. Note that
the closed form $\arcsin(\Rmin/\Rmax)-\delta$, obtained from the discriminant of \eqref{eq:c1}
alone, is only an \emph{upper bound} on $\psi_{\max}$.
\end{theorem}
\begin{proof}
For fixed $\beta$ the function $r\mapsto d^2+r^2-2dr\cos(\psi-\beta)$ is convex, so its maximum
over $[r_{\min},\Rmax]$ is attained at an endpoint. Maximizing over $\beta$ then requires
minimizing $\cos(\psi-\beta)$, attained at the endpoint of $[-\delta,\delta]$ farthest from
$\psi$, which gives $|\psi-\beta|=|\psi|+\delta=\Theta$ and yields \eqref{eq:maxdist}. Requiring
the maximum to be at most $\Rmin$ gives \eqref{eq:c1}--\eqref{eq:c2}. The discriminant of
\eqref{eq:c1} is non-negative iff $\cos^2\Theta\ge1-(\Rmin/\Rmax)^2$, i.e.
$|\psi|\le\arcsin(\Rmin/\Rmax)-\delta$ on the small-angle branch.
\end{proof}

For the numerical values of this work ($\Rmin=1000$~m, $\Rmax=1500$~m, $\delta=1^\circ$)
both constraints are active. At $\psi=0$ the admissible interval is
$d\in[500.1,1005.0]$~m: the lower end comes from the intersection angle and the upper end from
\eqref{eq:c2}, not from \eqref{eq:c1}, whose own upper root is $2499.4$~m. As $|\psi|$ grows the
interval defined by \eqref{eq:c1} shrinks while the upper root of \eqref{eq:c2} moves only
slightly, and the two cease to intersect at
\begin{equation}
\psi_{\max}=40.44^\circ,\qquad d=1003.7\ \text{m},
\label{eq:psimax}
\end{equation}
where the admissible interval degenerates to a single point. Substituting the largest admissible
radius gives $\arcsin(\Rmin/\Rmax)-\delta=40.81^\circ$, a strict upper bound that is never
attained. The practical consequences are that the second measurement must be placed \emph{off} the
first bearing by at most about $40^\circ$, and that its distance is bounded both from below (a
usable intersection angle) and from above (the emitter may be close to $S_1$).


\begin{remark}[The distance criterion is not sufficient for directional emitters]
\label{rem:dirsecond}
The condition $\max_{X\in\bF_1}\|S_2-X\|\le\Rmin$ guarantees reception only when reception depends
on distance alone. For a directional emitter it does not, and no distance bound can repair this: a
site may be arbitrarily close to the emitter and still lie outside the illuminated half-plane.
A concrete instance in the geometry of this paper is $S_1=(0,0)$, $S_2=(600,200)$ and an emitter at
$G=(900,0)$ with half-plane $H=\{(x,y):x+2y\le900\}$. The first measurement at $S_1$ returns a
bearing, and $\|S_2-G\|=361$~m is well inside $\Rmin=1000$~m, yet
$600+2\cdot200=1000>900$, so $S_2\notin H$ and no signal is received at the second site. With
$\delta=1^\circ$ the worst-case candidate distance from $S_2$ is $928$~m, so this pair passes
\eqref{eq:c1}--\eqref{eq:c2} comfortably: the test accepts a placement that does not work.

Theorem~\ref{thm:second} must therefore be read either as a statement about omni-directional
emitters, or as a statement about sites already known to lie in the illuminated half-plane. For
directional emitters the distance condition is necessary but not sufficient, and measurement
placement must instead be driven by the enclosure condition \eqref{eq:conv}; this is exactly the
role played by the census network of Section~\ref{sec:network}.
\end{remark}


\begin{remark}[Co-channel emitters]
\label{rem:cochannel}
Assumption (A3) excludes two transmitters occupying the same channel. Relaxing it does not
invalidate the geometry of this paper, but it does change what the bearing measurement is.
With $m\ge2$ co-channel emitters the sampled covariance acquires $m$ significant eigenvalues, the
MUSIC spectrum of Section~\ref{sec:rf} exhibits $m$ peaks, and a single reported bearing is no
longer an estimator of one emitter position but a sample from a mixture. The correct object is
then a \emph{set} of bearings per measurement, and the natural extension is to cluster the
bearings into $m$ tracks and run one evidence set per track. Because each track carries its own
enclosure condition, the network guarantee of Theorem~\ref{thm:detect} transfers unchanged to
each track; what has to be re-derived is the bookkeeping, not the geometry. Two practical
consequences are worth noting: the number of sources must be estimated rather than assumed, which
is standard in array processing but introduces its own error mode when two emitters are closely
spaced; and the absence certificate must be issued per track, so a channel is certified only when
every track has been resolved. We regard co-channel operation as an orthogonal extension of the
present model rather than a special case of it.
\end{remark}

\subsection{Detectability of orientation-unknown emitters}
\label{sec:detect}

We now state the central result. Let $\bP=\{P_1,\dots,P_m\}$ be a set of measurement locations
(a \emph{network}) and, for a candidate position $G$, let
$S_G=\{P\in\bP:\|P-G\|\le\Rmin\}$ be the network points within the worst-case reception radius.

For a candidate position $G$ and a reception radius $\rho>0$, let
\begin{equation}
S_G(\rho)=\{P\in\bP:\ \|P-G\|\le\rho\}
\label{eq:SGrho}
\end{equation}
be the set of network points that an emitter at $G$ with effective radius $\rho$ could reach. The
radius is unknown but bounded, $\rho\in[\Rmin,\Rmax]$ (Section~\ref{sec:problem}), and the
quantifier below is stated accordingly.

\begin{theorem}[Detectability at a fixed reception radius]
\label{thm:detect}
Fix $\rho>0$ and write $S_G=S_G(\rho)$. The network detects an emitter at $G$ for \emph{every}
orientation $d$ if and only if
\begin{equation}
G\in\conv(S_G(\rho)).
\label{eq:convrho}
\end{equation}
\end{theorem}
\begin{proof}
($\Leftarrow$) Let $d$ be arbitrary and let $H=\{X:(X-G)\cdot u(d)\ge0\}$ be the illuminated
half-plane. If no point of $S_G$ lay in $H$, then all of $S_G$ would lie in the open complementary
half-plane, whose closure is convex and hence contains $\conv(S_G)$; this would contradict
$G\in\conv(S_G)$ because $G$ lies on the boundary of $H$. Thus some $P\in S_G\cap H$ exists, and
$\|P-G\|\le\rho$, so the emitter is heard from $P$ regardless of $d$.
($\Rightarrow$) Suppose $G\notin\conv(S_G)$. By the separating hyperplane theorem there is a
direction $d$ with $(P-G)\cdot u(d)<0$ for all $P\in S_G$, i.e. every network point within the
reception radius lies strictly behind the emitter. Choosing this $d$ as the orientation makes the
emitter invisible from every point of $S_G$, and points outside $S_G$ are, by definition of
$S_G(\rho)$, beyond the reception radius $\rho$. Hence the emitter is not detected, a
contradiction.
\end{proof}

\begin{corollary}[Uniform guarantee over the admissible radii]
\label{cor:uniform}
The network detects an emitter at $G$ for every orientation $d$ \emph{and every} admissible
reception radius $\rho\in[\Rmin,\Rmax]$ if and only if
\begin{equation}
G\in\conv(S_G(\Rmin)).
\label{eq:conv}
\end{equation}
\end{corollary}
\begin{proof}
Since $\rho\mapsto S_G(\rho)$ is non-decreasing, $\conv(S_G(\Rmin))\subseteq\conv(S_G(\rho))$ for
all $\rho\ge\Rmin$, so \eqref{eq:conv} implies \eqref{eq:convrho} at every admissible radius by
Theorem~\ref{thm:detect}. Conversely, if $G\notin\conv(S_G(\Rmin))$ then the separating direction
constructed in the necessity part of that proof, applied at $\rho=\Rmin$, renders the emitter
invisible at an admissible radius.
\end{proof}

\begin{remark}[Which form is used where]
\label{rem:quantifier}
Equation~\eqref{eq:convrho} is the exact statement for a \emph{known} radius and
\eqref{eq:conv} is its worst-case counterpart for an \emph{unknown} one; the two coincide when
$\rho$ is known to equal $\Rmin$. Because $\rho$ is unknown in this paper we use the worst-case
form throughout, writing $S_G$ for $S_G(\Rmin)$. The distinction is not merely pedantic in the
other direction: for a known $\rho>\Rmin$ the enclosure condition \eqref{eq:convrho} is strictly
easier to satisfy than \eqref{eq:conv}, so the worst-case form is conservative, and
Section~\ref{sec:reliability} quantifies the reliability that this conservatism buys and costs.
\end{remark}

\begin{corollary}[The network must leave the region]
\label{cor:outside}
If the emitter can lie on the boundary of $\bA$, then any network that satisfies
\eqref{eq:conv} for all $G\in\bA,\ d$ must contain points outside $\bA$.
\end{corollary}
\begin{proof}
Take $G$ on the boundary with outward normal $n$ and consider $d=n$. For every $P\in\bA$ we have
$(P-G)\cdot n\le0$, with equality only if $P=G$. Hence no point of $\bA$ lies in the open
half-plane $H(G,n)^\circ$, so $S_G$ contains no point of $H(G,n)^\circ$ and
$G\notin\conv(S_G)$ whenever $G\notin\bP$.
\end{proof}

\begin{corollary}[Spacing]
\label{cor:spacing}
If \eqref{eq:conv} holds for all $G$ in a disk of radius $R$, then the largest empty circle of
$\bP$ has radius at most $\Rmin$, i.e. the covering radius of $\bP$ does not exceed $\Rmin$.
\end{corollary}
\begin{proof}
If some $G$ had $\|G-P\|>\Rmin$ for all $P$, then $S_G=\varnothing$ and $G\notin\conv(\varnothing)$,
violating \eqref{eq:conv}.
\end{proof}

\begin{figure*}[t]
\centering
\includegraphics[width=\textwidth]{en_net.pdf}
\caption{Detectability for emitters of unknown orientation. Panels (a) and (b) colour the region by whether \eqref{eq:conv} holds; the seven-point ring network that covers the region adequately for omni-directional emitters satisfies it on only about $24\%$ of the area, whereas the two-layer network satisfies it on all but a residual of measure tending to zero. The fraction is estimated by point sampling ($23.6\%$ on a $20$~m grid) and by the cell decomposition of Section~\ref{sec:network}, which converges to the same value as the mesh is refined. Panel (c) illustrates the two cases of Theorem~\ref{thm:detect} (the detectability condition):
for $G\notin\conv(S_G)$ there is an orientation (red arrow) that illuminates no network point.}
\label{fig:net}
\end{figure*}

Corollary~\ref{cor:spacing} shows that the classical covering condition is \emph{necessary} but,
by Theorem~\ref{thm:detect}, not sufficient: covering only ensures $S_G\ne\varnothing$, whereas
detection for all orientations requires that $S_G$ surrounds $G$. Fig.~\ref{fig:net} makes the gap
concrete: the seven-point ring network that covers the region adequately for omni-directional
emitters satisfies \eqref{eq:conv} on only about $24\%$ of the area, while the two-layer network satisfies it on all but a residual of measure tending to zero. This is the formal reason why
the dense lattice of Section~\ref{sec:network} is needed, and why a sparse covering network that
is perfectly adequate for omni-directional emitters fails badly when orientations are unknown.

\begin{remark}
Theorem~\ref{thm:detect} is tight and testable. In our implementation the condition is evaluated
numerically by computing the angular gaps of $S_G$ as seen from $G$: $G\in\conv(S_G)$ holds iff
some point of $S_G$ coincides with $G$ or the largest angular gap does not exceed $\pi$. We
cross-validated this criterion against an independent triangle-barycentre test on $10^5$ sample
points with zero discrepancies.
\end{remark}

\subsection{Certified absence and provable termination}

\begin{theorem}[Absence certificate]
\label{thm:cert}
Suppose $\bP$ satisfies \eqref{eq:conv} for every $G\in\bA$. If channel $c$ has been measured at
every point of $\bP$ and every reading was \texttt{no\_signal}, then no emitter occupies channel
$c$.
\end{theorem}
\begin{proof}
Assume an emitter exists at $G\in\bA$ with orientation $d$ and reception radius $\rho\ge\Rmin$.
By Corollary~\ref{cor:uniform} --- equivalently Theorem~\ref{thm:detect} applied at $\rho=\Rmin$ ---
there is $P\in\bP$ with $\|P-G\|\le\Rmin\le\rho$ and $P\in H(G,d)$, so the reading at $P$ would not have
been \texttt{no\_signal}; contradiction.
\end{proof}

\begin{corollary}[Decidable termination]
\label{cor:terminate}
Let $\bP$ satisfy \eqref{eq:conv}. A mission may terminate with the guarantee that every emitter
has been neutralized as soon as, for every channel, either a neutralization succeeded or the
channel has been measured at all points of $\bP$ with \texttt{no\_signal} readings.
\end{corollary}

Corollary~\ref{cor:terminate} is the operational meaning of certification. Without
Theorem~\ref{thm:cert} the platform can never distinguish ``all emitters neutralized'' from
``some emitter not yet found'', and must keep searching until an external timeout --- a
distinction whose cost we quantify in Section~\ref{sec:ablation}.


\subsection{Robustness to outliers: $q$-relaxed evidence sets}
\label{sec:robust}

The developments so far assume that every reading is truthful. In a real radio channel this
assumption fails: multipath reflections and antenna sidelobes produce \emph{spurious bearings}
--- a \texttt{direction} reading for a channel that is not present, or a bearing far from the true
one --- while shadowing and fading produce \emph{missed detections}. The two failure modes are not
symmetric, and neither is their effect on the certificate.

\begin{definition}[$q$-relaxed evidence set]
\label{def:qrelax}
For $q\in\mathbb{N}_0$, the set $E_c^{(q)}$ consists of those $X\in\bA$ that violate at most
$q$ of the constraints accumulated for channel $c$.
\end{definition}

The rule $E\leftarrow E\cap C$ of Table~\ref{tab:ops} is the case $q=0$. For $q>0$ the update is
implemented by maintaining a violation counter $\nu(X)$ per cell and setting
$E^{(q)}=\{X:\nu(X)\le q\}$; this adds no asymptotic cost and no extra memory beyond one small
integer per cell.

\begin{theorem}[Robust clearance]
\label{thm:robustclear}
Assume that at most $q$ of the readings taken for channel $c$ are spurious. If
$\MEC(E_c^{(q)})\le r_{\mathrm{clr}}-\eta$, then a single neutralization attempt at the centre of
$E_c^{(q)}$ succeeds with certainty.
\end{theorem}
\begin{proof}
A truthful constraint is satisfied by the true emitter position, so $G_c$ violates at most the $q$
spurious constraints and hence $G_c\in E_c^{(q)}$. The argument of Theorem~\ref{thm:clear} then
applies verbatim with $E_c$ replaced by $E_c^{(q)}$.
\end{proof}

Absence certification needs separate treatment. A spurious \emph{bearing} is harmless for it,
because spurious bearings can only add detections and therefore cannot hide the detection that
Theorem~\ref{thm:detect} guarantees; this is why the certificate of Theorem~\ref{thm:cert} is
untouched. A \emph{missed detection}, however, can remove exactly the detection that the
certificate relies on, and it can also shrink the raster to empty and thereby produce a
\emph{false absence certificate}. Two consequences follow.

The next theorem concerns the first of these two mechanisms only, and we make the meaning of $q$
precise before stating it. \emph{Missed detections are treated separately in
Section~\ref{sec:burst} and are not covered by the redundancy condition below.} The two failure
modes enter the certificate through opposite mechanisms. A spurious \emph{positive} detection
inflates the detection count of an otherwise empty channel, and is therefore defeated by
redundancy: enough truthful detections out-vote the spurious ones. A missed detection, by
contrast, removes a detection on which Theorem~\ref{thm:cert} depends, and no amount of network
redundancy repairs it --- only repeated measurement in time can, at the cost of the probabilistic
guarantee \eqref{eq:rounds}.

\begin{theorem}[Redundancy condition for robust absence]
\label{thm:redundancy}
Let the network satisfy \eqref{eq:conv} and suppose that at most $q$ detections per channel are
spurious. If in addition the network has \emph{detection redundancy}
\begin{equation}
r=\min_{G\in\bA,\ d}\ \bigl|\{P\in\bP:\|P-G\|\le\Rmin,\ P\in H(G,d)\}\bigr|\ >\ q,
\label{eq:redundancy}
\end{equation}
then a channel that has been measured at every point of $\bP$ and for which at most $q$
detections were recorded is certified absent.
\end{theorem}
\begin{proof}
If an emitter of orientation $d$ occupies channel $c$, \eqref{eq:redundancy} guarantees at least
$r$ network points from which it is truthfully heard, so the number of recorded detections is at
least $r>q$, contradicting the hypothesis.
\end{proof}

\begin{proposition}[Redundancy of the proposed networks]
\label{prop:r1}
For the disk geometry with $R=1800$~m and $\Rmin=1000$~m, the seven-point ring network has $r=0$
and the 25-point two-layer network has $r=1$, the minimum being attained at boundary positions
where every network point within $\Rmin$ lies on the interior side and the outward half-plane
therefore contains only points of the outer ring.
\end{proposition}

Proposition~\ref{prop:r1} has a practical reading that we consider important:
\emph{the detectability condition \eqref{eq:conv} is what makes absence certifiable, but it is not
what makes absence certifiable \textbf{robustly}}. Robustness to $q$ spurious readings requires
$r\ge q+1$, a strictly stronger network property. In the disk geometry this is expensive: across
the families we examined (an additional outer ring, lattices down to $700$~m spacing, and
extended inner radii, up to $49$ points) the value $r=1$ persisted, because a boundary position
always admits a half-plane containing only the single outermost point. Constructing a
redundancy-$r$ network of minimum tour length is therefore an open design problem, which we state
in Section~\ref{sec:discussion}.

\begin{remark}[Implementation]
When $q>0$ the implementation (i)~uses the $q$-relaxed intersection of
Definition~\ref{def:qrelax}, (ii)~disables the raster-emptying path to absence, and
(iii)~disables the distance inference from \texttt{no\_signal} for omni-directional emitters,
because a missed detection would otherwise trim the true position. Only the census certificate of
Theorem~\ref{thm:cert} is then accepted as a proof of absence. Sections~\ref{sec:outlier} and
\ref{sec:discussion} quantify what this buys and what it costs.
\end{remark}

\subsection{Lower bounds for network design}
\label{sec:bounds}

Let $k=|\bP|$ and let $L(\bP)$ be the length of the shortest tour from the origin visiting all
points of $\bP$.

\begin{theorem}[Point-count lower bound]
\label{thm:klb}
For an omni-directional emitter with reception radius at least $\Rmin$ that may lie anywhere in
$\bA$, any network satisfying the coverage requirement $\bA\subseteq\bigcup_i B(P_i,\Rmin)$ has
\begin{equation}
k\ \ge\ \left\lceil\frac{\pi}{\varphi_{\max}}\right\rceil=6,
\qquad
\varphi_{\max}=\arccos\!\sqrt{1-\frac{\Rmin^2}{R^2}}=\arcsin\!\frac{\Rmin}{R}=33.7^\circ .
\label{eq:klb}
\end{equation}
\end{theorem}
\begin{proof}
The circle $\partial\bA$ must be covered. A point $P$ at distance $\rho$ from the origin whose
disk of radius $\Rmin$ intersects $\partial\bA$ covers an arc of half-angle $\varphi(\rho)$ with
$\cos\varphi(\rho)=(\rho^2+R^2-\Rmin^2)/(2\rho R)$. Differentiating, $\varphi$ is maximized at
$\rho^\ast=\sqrt{R^2-\Rmin^2}$, giving
$\cos\varphi_{\max}=(\rho^{\ast2}+R^2-\Rmin^2)/(2\rho^\ast R)=\sqrt{1-\Rmin^2/R^2}$, i.e.
$\varphi_{\max}=\arcsin(\Rmin/R)$. For $R=1800$~m, $\Rmin=1000$~m:
$\varphi_{\max}=33.75^\circ$, so each point covers at most $67.5^\circ$ of the boundary and
$k\ge\lceil360/67.5\rceil=6$.
\end{proof}

\begin{theorem}[Route-length lower bound]
\label{thm:llen}
Any tour from the origin that visits all points of a network covering $\bA$ has
\begin{equation}
L(\bP)\ \ge\ \frac{\pi(R^2-\Rmin^2)}{2\Rmin}\ =\ 3519\ \text{m}
\quad (R=1800,\ \Rmin=1000).
\label{eq:llen}
\end{equation}
\end{theorem}
\begin{proof}
A path of length $L$ sweeps at most $2\Rmin L+\pi\Rmin^2$ units of area with disks of radius
$\Rmin$. Covering $\bA$ of area $\pi R^2$ therefore requires
$2\Rmin L+\pi\Rmin^2\ge\pi R^2$.
\end{proof}

Theorem~\ref{thm:klb} bounds the number of points from below by six. Whether six points actually
suffice is a question about the covering number of the disk: writing $\varrho_n$ for the smallest
radius with which $n$ equal circles can cover a unit disk, six points suffice if and only if
$\varrho_6\le\Rmin/R=0.55556$. Two independent lines of evidence indicate that this fails.
First, local searches over the natural families (six on a ring, one centre plus five, and
asymmetric starts), with the covering radius evaluated \emph{exactly} from the Voronoi vertices
and the boundary circle rather than from a sample grid, found \emph{no} six-centre configuration
with covering radius below $0.5774$. This is numerical evidence about the configurations we
searched and is \emph{not} a lower bound on $\varrho_6$; it is reported only for orientation.
Second, the optimal covering radius for six equal circles is
$\varrho_6=0.555905\ldots$~\cite{bezdek1983,nica2011},
which still exceeds the threshold. Six points are therefore insufficient, seven are necessary,
and the layout of
Section~\ref{sec:network} attains the covering requirement with a margin of $3.1\%$
($969$~m against $\Rmin=1000$~m).

We flag two points of book-keeping. The conclusion rests on the covering-number literature and
not on an enumeration of configurations by the authors; and the tabulated margin is thin
($0.06\%$), so a reader who distrusts that value should read our contribution as ``seven points
suffice for this geometry'' rather than ``seven points are optimal''. A direct, assumption-free
proof that six points are insufficient for $\Rmin/R=0.55556$ would be welcome and is left open.


%=====================================================================
\section{Method}
\label{sec:method}

\subsection{Overview}
The method has three layers. \textbf{(L1) A census network} makes one measurement on every
unresolved channel at each network point, which detects every emitter whenever the enclosure
condition \eqref{eq:conv} holds (Theorem~\ref{thm:detect}) and provides the evidence needed for the
absence certificate
(Theorem~\ref{thm:cert}). \textbf{(L2) Evidence-driven approach} moves the platform towards the
centre of the current evidence set of a chosen channel, re-measuring on arrival. \textbf{(L3)
Certified clearance} neutralizes a channel as soon as Theorem~\ref{thm:clear} applies, and falls
back to a provably complete raster sweep when the evidence set cannot be shrunk further.

\subsection{Network design}
\label{sec:network}

For omni-directional emitters, Corollary~\ref{cor:spacing} requires the covering radius to be at
most $\Rmin$; for unknown orientations Theorem~\ref{thm:detect} additionally requires
\emph{enclosure}, and Corollary~\ref{cor:outside} forces part of the network outside $\bA$. These
requirements suggest the following optimization:
\begin{equation}
\min_{d,\,R_{\mathrm{in}},\,n_{\mathrm{out}},\,\rho_{\mathrm{out}}}\ L(\bP)
\quad\text{s.t.}\quad \forall G\in\bA:\ G\in\conv(S_G),
\quad
\bP=\operatorname{hex}(d,R_{\mathrm{in}})\cup\operatorname{ring}(n_{\mathrm{out}},\rho_{\mathrm{out}}),
\label{eq:netopt}
\end{equation}
i.e. a dense inner hexagonal lattice of nearest-neighbour spacing $d$ covering a disk of radius
$R_{\mathrm{in}}$, plus an outer ring of $n_{\mathrm{out}}$ points at radius
$\rho_{\mathrm{out}}$. Problem \eqref{eq:netopt} is a combinatorial optimization over four
parameters with a convex inner problem, which we solve by exhaustive search on a grid followed by
local refinement~\cite{boyd2004}.

The intuition --- which the experiments confirm --- is that the outer ring need only be placed
\emph{just} outside the region: it exists to enclose boundary emitters that radiate outwards. This
is markedly more efficient than extending a single lattice far beyond the boundary. Solving
\eqref{eq:netopt} over the four parameters yields
$d=1000$~m, $R_{\mathrm{in}}=1800$~m ($13$ points), $n_{\mathrm{out}}=12$,
$\rho_{\mathrm{out}}=1900$~m, i.e. $k=25$ points and a route of $17.7$~km, which is $42\%$ shorter
than the single-layer lattice satisfying the same enclosure condition ($31$ points, $30.7$~km). We do
not rely on a single finite-grid check. The condition \eqref{eq:conv} is certified analytically on
all but a shrinking residual, and that residual is then evaluated numerically. A finite grid check
alone cannot establish \eqref{eq:conv} for every point of $\bA$, so we use a \emph{conservative
cell certificate}. For a square cell $C$ with circumradius $\eta$ define
\begin{equation}
S_C^-=\{P\in\bP:\ \textstyle\max_{G\in C}\|P-G\|\le\Rmin\}.
\label{eq:Sminus}
\end{equation}
Since $\|P-\cdot\|$ is convex the maximum is attained at a corner of $C$, so $S_C^-$ is decided
exactly. Every $P\in S_C^-$ lies within $\Rmin$ of \emph{every} $G\in C$, hence
$S_C^-\subseteq S_G$ for all $G\in C$; and because $\conv(S_C^-)$ is convex while $C$ is the
convex hull of its four corners, it suffices to test the corners:
\begin{equation}
C\subseteq\conv(S_C^-)\quad\Longrightarrow\quad G\in\conv(S_G)\ \ \text{for all }G\in C.
\label{eq:cellcert}
\end{equation}
Note that \eqref{eq:cellcert} does not require the absence of a circle
$\partial B(P_i,\Rmin)$ through $C$; cells that fail the test are simply subdivided.

Applied to the $25$-point network with adaptive subdivision, \eqref{eq:cellcert} certifies the
arena except for a residual of $32.8$~m$^2$ at an effective resolution of $0.63$~m and
$9.4$~m$^2$ at $0.31$~m, against a total area of $1.018\times10^{7}$~m$^2$ --- a fraction of about
$3\times10^{-6}$. \emph{Every} residual cell intersects some circle $\partial B(P_i,\Rmin)$; the
uncertified set is therefore a shrinking neighbourhood of the union of $m$ circles, whose area
tends to zero with the mesh. This residual is a limitation of the conservative certificate rather
than evidence of a gap in the network: on it we evaluate \eqref{eq:conv} directly on a
$21\times21$ sub-grid per cell, no point fails, and the largest angular gap found is
$165^\circ$ against the $180^\circ$ limit, a margin of $15^\circ$. We therefore state the
guarantee as: \emph{enclosure is proved on all of $\bA$ except a set whose measure tends to zero,
and verified numerically on the remainder with a $15^\circ$ margin.} Closing the residual by
partitioning along the circle arcs rather than by axis-aligned subdivision is left to future
work. We make the logical status explicit, because the paper's
objective is stated for every realization. The results of Section~\ref{sec:theory} are
unconditional: they apply to \emph{any} network that satisfies \eqref{eq:conv} at every point of
$\bA$. What is established for the specific $25$-point construction is that it satisfies
\eqref{eq:conv} analytically on all of $\bA$ except a residual of measure tending to zero, and
numerically --- with a $15^\circ$ margin --- on that residual. Consequently the mission-level
guarantees of Theorems~\ref{thm:cert} and \ref{cor:terminate} are conditional, for this particular
network, on that verification; a reader who requires the premise to be established purely
analytically should read the $25$-point results as conditional on it. The weaker
covering requirement needed in the omni-directional scenario is certified by the conservative test
$\min_i\|x-P_i\|\le\Rmin-\eta$, $\eta=h\sqrt2/2$: at $h=5$~m all $405\,605$ cells pass with zero
failures. The seven-point network certifies only about $24\%$ of the area against the stronger enclosure condition, which is precisely why it appears only in the omni-directional scenario. We note that this fraction is resolution-dependent: the cell decomposition reports $28.0\%$, $25.6\%$ and $24.5\%$ at $h=20$, $10$ and $5$~m, so it should be read as a convergent estimate rather than a mesh-specific number. For the two-layer network, direct pointwise evaluation reports $1.0000$ at every tested resolution; that numerical result is distinct from the conservative cell certificate above and should not be conflated with it.

A counter-intuitive finding is that the spacing requirement is \emph{not} the one suggested by the
Delaunay geometry. Because the circumradius of the triangles of a hexagonal lattice is
$d/\sqrt3$, one might expect $d\le\sqrt3\Rmin\approx1732$~m to suffice in the interior. Numerically
$d=1200$~m yields only $65.5\%$ of positions with the property \eqref{eq:conv}; the failures are
concentrated in an annulus of radius $200$--$600$~m around each lattice point, where exactly one
network point lies within $\Rmin$ and therefore cannot enclose $G$. Hence $d\le\Rmin$ is required.

For omni-directional emitters the census network reduces to the seven-point layout of
Section~\ref{sec:bounds}: a centre point plus a six-point ring at $\rho=1200$~m, whose covering
radius is $969$~m and whose route is $7.2$~km.

\subsection{Evidence update and clearance}
Algorithm~\ref{alg:update} summarizes the evidence update. Two representations are maintained in
parallel: a $10$~m raster, which can express the non-convex constraints induced by
\texttt{no\_signal} and by failed neutralizations, and an exact convex polygon obtained by
intersecting bearing wedges, which is tighter for the clearance test. Both are outer
approximations, so the more permissive of the two tests is used.

\begin{algorithm}[t]
\caption{Evidence update for channel $c$}
\label{alg:update}
\begin{algorithmic}[1]
\State \textbf{input} reading $z$, position $p$, channel $c$, current set $E_c$
\If{$z$ is a bearing $\bar\theta$}
    \State $E_c \gets E_c\cap W(p,\bar\theta)\cap \bar B(p,\Rmax)$
\ElsIf{$z$ is \texttt{near}}
    \State $E_c \gets E_c\cap \bar B(p,r_{\mathrm{near}})$
\ElsIf{$z$ is \texttt{no\_signal} and channel $c$ is omni-directional}
    \State $E_c \gets E_c\setminus B(p,\Rmin)$
\ElsIf{$z$ is a failed neutralization}
    \State $E_c \gets E_c\setminus B(p,r_{\mathrm{clr}})$
\EndIf
\State \Return $E_c$
\end{algorithmic}
\end{algorithm}

\subsection{Rolling-horizon decision}
At each stop the platform (i) measures every unresolved channel in ascending order, which also
minimizes channel switching, (ii) neutralizes any channel for which
Theorem~\ref{thm:clear} applies, and (iii) otherwise selects a target as follows.

\begin{itemize}
\item If some channel has a \emph{clearable} evidence set, move to its MEC centre and neutralize.
\item Else, if some channel has bearing evidence, step towards the centroid of its evidence set,
with a step length capped at $400$~m; re-measure on arrival.
\item Else, if the census is incomplete, proceed to the next census point.
\end{itemize}

Two fallbacks handle the directional case, in which \texttt{no\_signal} conveys no information and
the evidence set can freeze. First, an \emph{adaptive probe sequence}: points are laid out inside
the evidence set with spacing $\operatorname{clip}(\max(1.2\,\ell_{\min},\ell_{\max}/8),40,400)$~m,
where $\ell_{\min},\ell_{\max}$ are the minor and major axes of the set obtained by principal
component analysis, and are visited until a bearing is recovered. Spacing proportional to the
bounding-box diagonal --- the obvious choice --- is inadequate because the set is frequently a thin
wedge of $1500\times50$~m for which such spacing produces only two or three probes. Second, a
\emph{provably complete raster sweep}: neutralization attempts on a grid of spacing
$r_{\mathrm{clr}}\sqrt2=28.3$~m cover the evidence set in the sense that every point of the set is
within $r_{\mathrm{clr}}$ of a grid point, so the sweep necessarily succeeds.

\subsection{Complexity}
Each update is $O(N_{\mathrm{cells}})$ with $N_{\mathrm{cells}}\approx10^5$ for a $10$~m raster,
implemented as vectorized array operations; the polygon test is $O(m\log m)$ in the number of
accumulated bearings. A mission performs $O(|\bP|\cdot N_c)$ measurements, i.e. a few hundred,
which is negligible compared with the travel time.

%=====================================================================
\section{Experiments}
\label{sec:experiments}

\subsection{Setup}
All experiments use a discrete-event simulator that reproduces the sensor and timing model of
Section~\ref{sec:problem} exactly; the simulator was validated against the reference
implementation of the application programming interface and reproduces its documented timing
example to the second. Sample sizes differ between experiments and are stated in every table caption: the headline
S1 result uses $200$ episodes, the head-to-head comparison with the baselines uses $30$ per method,
the ablation uses $30$ (S1) and $15$ (S2), the generalization sweeps and the stress tests use $20$,
and the non-ideal-channel study uses $12$. Emitters are drawn uniformly in $\bA$, channels are
drawn without replacement from $\{1,\dots,20\}$, reception radii are uniform in $[1000,1500]$~m,
and orientations are uniform in $[0,360^\circ)$. Two scenarios are studied: \textbf{S1}
(omni-directional) and \textbf{S2} (directional), the latter with all emitters directional unless
stated otherwise. A mission is \emph{certified} if every channel is either neutralized or proved
empty by Theorem~\ref{thm:cert}; missions are cut off at a virtual-time budget of $60000$~s.

All statistics report paired tests between methods evaluated on identical seeds: an exact
two-sided sign test, together with a $95\%$ bootstrap confidence interval ($2000$ resamples) of
the mean difference.

\subsection{Baselines}
\label{sec:maincmp}
Comparing against weak baselines would be uninformative, so all baselines are given exactly the
same information as the proposed method --- the bearing readings, the \texttt{near} and
\texttt{no\_signal} outcomes and the neutralization results --- and are denied only the
bounded-error evidence-set machinery that is our contribution. The four baselines are:
\textbf{Lawnmower}, a complete boustrophedon sweep with $1800$~m row spacing and $300$~m stop
spacing, which homes on any channel it hears; \textbf{RandomWaypoints}, the same homing logic
driven by $40$ random waypoints; \textbf{HomingOnly}, no census at all, which simply chases
whichever channel it currently hears; and \textbf{GridClear}, a provably complete raster sweep of
neutralization attempts at $28.3$~m spacing that uses no bearing information whatsoever and is
reported analytically because its cost is two orders of magnitude larger. A sweep of the
$1.018\times10^7$~m$^2$ region at this spacing visits $12\,727$ points per channel, costing
$3$~s each, and requires a boustrophedon path of $359\,940$~m; the total is
$110\,170$~s per channel and $2.2\times10^6$~s for the $20$ channels.

A detail worth recording is that a naive homing law with a fixed step oscillates. With a $500$~m
step and an emitter $660$~m away, the platform repeatedly overshoots and returns, clearing
nothing ($0.8\%$ success in our first implementation). All baselines therefore use an adaptive
step that is halved whenever the bearing reverses by more than $120^\circ$; this alone raises the
\texttt{HomingOnly} success rate from $0.8\%$ to $31\%$. Reporting the fixed-step variant would
have constituted a straw man.

\noindent\textbf{Scope of the comparison.} One family of methods is deliberately not run as a
baseline: Bayesian and information-theoretic active search, from
infotaxis~\cite{vergassola2007} to learning-based adaptive informative path planning for
multi-robot target mapping~\cite{popovic2024,vashisth2026}. Those methods optimise an expected information gain or detection
probability under a \emph{prior} over emitter locations and a \emph{likelihood} model of the
field, and they terminate on a budget or on a confidence threshold. Neither input exists here:
the sensor returns a bounded-error bearing, or the absence of a signal, with no field strength,
no range and no prior on the number of emitters, and the mission must terminate on a
\emph{proof} rather than on a posterior. Running such a method would therefore compare
termination criteria rather than search efficiency, and the $4.8$-fold effect we measure is
precisely the value of the certificate \emph{within} our own formulation. The baselines we do
run are representative non-certifying strategies operating under the same information, and
\texttt{GridClear} supplies the certifying but information-free extreme.

\begin{table*}[t]
\caption{Comparison against baselines on identical seeds ($30$ episodes per method per
scenario). All methods receive the same virtual-time budget of $60\,000$~s, which is $12.9$ times
the time the proposed method needs in S1 and $6.8$ times in S2. \texttt{GridClear} is analytic
rather than a budgeted run; its entries are seconds per channel, in both scenarios.}
\label{tab:main}
\centering
\small
\setlength{\tabcolsep}{5pt}
\begin{tabular}{lcccccc}
\toprule
& \multicolumn{3}{c}{Scenario S1 (omni-directional)} & \multicolumn{3}{c}{Scenario S2 (directional)} \\
\cmidrule(lr){2-4}\cmidrule(lr){5-7}
Method & Cleared & Certified & Total time (s) & Cleared & Certified & Total time (s) \\
\midrule
\textbf{Proposed}   & \textbf{1.000} & \textbf{1.000} & \textbf{4637} & \textbf{1.000} & \textbf{1.000} & \textbf{8787} \\
Lawnmower           & 0.509 & 0.000 & 35799 & 0.227 & 0.000 & 26261 \\
RandomWaypoints     & 0.513 & 0.067 & 51993 & 0.255 & 0.000 & 54471 \\
HomingOnly          & 0.308 & 0.000 & 55055 & 0.109 & 0.000 & 51066 \\
\midrule
GridClear (analytic) & 1.000 & --- & $1.10\times10^5$ & 1.000 & --- & same \\
\bottomrule
\end{tabular}
\end{table*}

Table~\ref{tab:main} reports the main comparison. The common budget of $60\,000$~s is $12.9$
times what the proposed method needs in S1 ($4637$~s) and $6.8$ times what it needs in S2
($8787$~s). Even so, the strongest baseline clears $51\%$ of the emitters in S1 and $25.5\%$ in S2.
Certification is essentially out of reach for the baselines: in S2 none of them certifies a single
mission, and in S1 only RandomWaypoints does, in $6.7\%$ of episodes. With only coverage-based
reasoning they cannot rule out an undetected emitter, so they run until the budget expires. The paired sign tests
reject equality in every case ($30/30$ wins, $p<2\times10^{-9}$); the mean differences in cleared
fraction are $+0.49$ (Lawnmower), $+0.49$ (RandomWaypoints) and $+0.69$ (HomingOnly) for S1, with
$95\%$ confidence intervals that exclude zero by a wide margin.

When the budget is reduced to $5000$~s, i.e. approximately the time the proposed method uses, the
gap widens: the proposed method clears $98.3\%$ of emitters and certifies $80\%$ of missions, while
the baselines clear $18.5\%$, $19.9\%$ and $25.4\%$.

The analytic entry for GridClear quantifies why bearing information is indispensable. A raster
sweep of the whole region at $28.3$~m spacing requires $12\,727$ neutralization attempts per
channel, i.e. $38\,181$~s of actions, plus $71\,988$~s of travel, for a total of
$110\,170$~s per channel; over $20$ channels this reaches $2.2\times10^6$~s, far beyond any
practical budget. The proposed method replaces this brute force by a few hundred measurements.

\begin{figure}[t]
\centering
\includegraphics[width=\columnwidth]{en_cmp.pdf}
\caption{Cleared fraction and certified-completion rate for the proposed method and the
baselines, under a common budget of $60\,000$~s --- $12.9$ times what the proposed method needs in
S1 and $6.8$ times in S2.}
\label{fig:cmp}
\end{figure}


\subsection{Isolating the planning contribution}
\label{sec:planning}

The baselines of Section~\ref{sec:maincmp} differ from the proposed method in every respect at
once, and the ablation of Section~\ref{sec:ablation} removes whole components. Neither answers the more basic
question: \emph{given the same evidence sets and the same certificate,
how much does the planning policy itself contribute?} Table~\ref{tab:planning} answers it by
holding the inference layer fixed and changing exactly one planning decision at a time.

\begin{table}[t]
\caption{Planning isolation. The evidence-set inference, the clearance criterion and the absence
certificate are identical in every row; only the planning decision shown changes. Figures are
means over $30$ (S1) and $20$ (S2) episodes. All rows clear $100\%$ of emitters and certify
$100\%$ of missions, so the differences are purely economic.}
\label{tab:planning}
\centering
\small
\begin{tabular}{lcc}
\toprule
Planning decision & S1 time (s) & S2 time (s) \\
\midrule
\textbf{Proposed} (NN$+$2-opt order, nearest channel first, centroid) & \textbf{4637} & \textbf{8887} \\
Census points in generation order & $4637$ & $11680$ \\
Census points in random order & $5672$ & $15775$ \\
Channels processed round-robin & $7219$ & $11504$ \\
Channels processed largest-set-first & $7192$ & $11597$ \\
Bounding-box centre instead of centroid & $4614$ & $8880$ \\
\bottomrule
\end{tabular}
\end{table}

Three observations follow. First, the tour matters: abandoning the Euclidean tour for a random
order costs $22\%$ in S1 and $78\%$ in S2, because the census accounts for roughly half of the
mission. (In S1 the generation order already coincides with the tour, so that row is a control and
changes nothing.) Second, the \emph{channel} order matters at least as much --- $+56\%$ in S1 and
$+30\%$ in S2 --- even though it changes no geometry at all: visiting nearby emitters in one pass
saves the repeated traversals that a round-robin policy incurs. Third, the choice of point estimate
inside the evidence set is essentially irrelevant ($-0.5\%$), which is reassuring: the method does
not depend on a lucky estimator, only on the set itself.

The separation is therefore clean. The inference layer determines \emph{whether} the
mission succeeds (ablation A1--A2: up to $-58\%$), whereas the planning layer determines
\emph{how fast} it succeeds (up to $+78\%$), and the two are independent.

\subsection{Ablation study}
\label{sec:ablation}

Table~\ref{tab:abl} removes one component of the method at a time. Four conclusions follow.

\begin{table}[t]
\caption{Ablation. S1: $30$ episodes; S2: $15$ episodes.}
\label{tab:abl}
\centering
\small
\begin{tabular}{lcccc}
\toprule
& \multicolumn{2}{c}{S1} & \multicolumn{2}{c}{S2} \\
\cmidrule(lr){2-3}\cmidrule(lr){4-5}
Variant & Cleared & Time (s) & Cleared & Time (s) \\
\midrule
Full                 & \textbf{1.000} & \textbf{4637} & \textbf{1.000} & \textbf{8871} \\
A1 no evidence set   & 0.420 & 4345 & 0.710 & 8309 \\
A2 no MEC criterion  & 1.000 & 5703 & 0.672 & 28966 \\
A3 no certificate    & 1.000 & 4637 & 1.000 & 42809 \\
A4 sparse network    & ---   & ---  & 0.996 & 47430 \\
A5 no probe sequence & 1.000 & 4646 & 0.995 & 8775 \\
A6 no raster sweep   & 1.000 & 4637 & 0.981 & 8895 \\
\bottomrule
\end{tabular}
\end{table}

\textbf{(i) Evidence-set inference is the performance pillar.} Replacing it by plain bearing
following (A1) collapses the cleared fraction from $1.000$ to $0.420$ in S1 and to $0.710$ in S2.

\textbf{(ii) The clearance criterion is equally important.} Without the MEC test (A2) the platform
repeatedly attempts neutralizations that cannot succeed; the cleared fraction in S2 falls to
$0.672$ and the mission time grows by a factor of $3.3$.

\textbf{(iii) The certificate is the most distinctive contribution.} Removing it (A3) leaves the
cleared fraction unchanged --- the same emitters are found --- but in S2 the mission time grows from
$8871$~s to $42809$~s, a factor of $4.8$, because the platform can no longer tell that it is
finished and searches until the budget is exhausted. In S1 the effect is null ($4637$~s with and
without), because there the raster path already resolves empty channels and the certificate is
redundant; the certificate matters exactly where \texttt{no\_signal} carries no distance
information. This quantifies the value of decidability
directly.

\textbf{(iv) The two fallbacks matter only for the last percent.} The probe sequence (A5) and the
raster sweep (A6) change the average cleared fraction by less than $2\%$, but they are exactly what
raises the success rate and the certification rate from $98.1\%/73.3\%$ to $100\%$.

\begin{figure*}[t]
\centering
\includegraphics[width=\textwidth]{en_gen.pdf}
\caption{Generalization. Each panel sweeps one axis; error bars are one standard deviation.
With the two-layer network (S2) every configuration attains a cleared fraction of $1.000$; the
green S1 curve in the directional panel is the single-layer ring network and is discussed in the
text.}
\label{fig:gen}
\end{figure*}


\subsection{Replacing the bearing model with an array-processing chain}
\label{sec:rf}

Every result so far assumes the abstract sensor model \eqref{eq:sensor}: a bearing equal to the
true angle plus a deterministic error bounded by $\delta=1^\circ$. That assumption is convenient
but untested, and in a real radio channel it is exactly what fails first. We therefore replace the
model with a complete direction-finding chain, simulated at sample level.

\noindent\textbf{The chain.} A narrowband far-field model is used. The platform carries a
uniform circular array of $M$ elements at radius $R_\lambda$ wavelengths. The channel from an
emitter at angle $\theta_0$ consists of a line-of-sight path plus $K$ specular reflections with
relative amplitude $\rho$, angles spread uniformly within $\pm\sigma$ and uniformly random phases.
Snapshots $x(t)=\sum_k\alpha_k\,a(\theta_k)s(t)+n(t)$ are formed with a constant-modulus pilot at
a prescribed SNR, the sample covariance is estimated from $N$ snapshots, and MUSIC assuming a
single source returns the strongest spectral peak. The receiver does not know $K$; consequently a
strong reflection can capture the peak and the reported bearing is simply wrong. The estimator is
deterministic for a given (position, channel) pair, so the structural property discussed in
Section~\ref{sec:problem} --- repeated measurements at one location do not help --- is preserved.

\noindent\textbf{Does the $\delta=1^\circ$ assumption survive?} Table~\ref{tab:rf} reports the error
distribution over $4000$ random geometries per configuration.

\begin{table*}[t]
\caption{Bearing error of the array chain ($4000$ trials each). The $\delta=1^\circ$ bounded-error
assumption of the analytical model holds only in the line-of-sight case.}
\label{tab:rf}
\centering
\small
\begin{tabular}{lcccccc}
\toprule
Channel & $p_{50}$ & $p_{90}$ & $p_{99}$ & $>5^\circ$ & $>20^\circ$ & $>45^\circ$ \\
\midrule
LOS only, SNR $20$~dB                & $0.09^\circ$ & $0.22^\circ$ & $0.34^\circ$ & $0.0\%$ & $0.0\%$ & $0.0\%$ \\
LOS only, SNR $12$~dB                & $0.22^\circ$ & $0.55^\circ$ & $0.85^\circ$ & $0.0\%$ & $0.0\%$ & $0.0\%$ \\
$1$ multipath, $\rho=0.6$            & $4.63^\circ$ & $12.2^\circ$ & $18.9^\circ$ & $47.3\%$ & $0.4\%$ & $0.0\%$ \\
$2$ multipath, $\rho=0.6$            & $6.90^\circ$ & $17.9^\circ$ & $52.1^\circ$ & $62.5\%$ & $7.4\%$ & $1.3\%$ \\
$3$ multipath, $\rho=0.8$            & $10.5^\circ$ & $30.5^\circ$ & $52.1^\circ$ & $75.0\%$ & $23.9\%$ & $2.9\%$ \\
\midrule
$2$ multipath, $M{=}16$, $R_\lambda{=}1.0$ & $2.77^\circ$ & $8.28^\circ$ & --- & --- & $3.8\%$ & $0.0\%$ \\
$2$ multipath, $\rho{=}0.3$, $M{=}16$, $R_\lambda{=}1.0$ & $1.22^\circ$ & $3.10^\circ$ & --- & --- & $0.0\%$ & $0.0\%$ \\
\bottomrule
\end{tabular}
\end{table*}

The bounded-error model is \emph{empirically justified} in a line-of-sight-dominated channel: at
$12$~dB the $99$th percentile error is $0.85^\circ$, comfortably inside $\delta=1^\circ$. With
multipath it collapses, and the failure is qualitative rather than quantitative: $7$--$24\%$ of
bearings are off by more than $20^\circ$ and a few percent by more than $45^\circ$. These are
precisely the spurious bearings that Section~\ref{sec:robust} treats analytically, now generated
by a physical model rather than postulated.

\noindent\textbf{What fixes it, and what does not.} Increasing the number of snapshots from $48$ to
$256$ barely helps ($p_{50}$ from $7.45^\circ$ to $6.98^\circ$): the error is dominated by the
multipath \emph{bias}, not by noise, so integration does not remove it. Increasing the array
\emph{aperture} does: a $16$-element array of radius $1.0\lambda$ reduces $p_{50}$ to
$2.77^\circ$ and eliminates the $>45^\circ$ tail entirely. Combining aperture with a moderate
reflection coefficient gives $p_{90}=3.1^\circ$ with no gross errors at all.

\begin{table}[t]
\caption{End-to-end runs with the array chain in place of the bearing model ($10$ episodes each,
scenario S2). The mission succeeds exactly when the sensor resolves the multipath.}
\label{tab:rfrun}
\centering
\small
\begin{tabular}{lccccc}
\toprule
Sensor / channel & $\delta$ & Cleared & \textbf{False abs.} & Certified & Time (s) \\
\midrule
$M{=}8$, $R_\lambda{=}0.35$, $\rho{=}0.6$ & $5^\circ$ & $0/10$ & $74$ & $0/10$ & 28243 \\
$M{=}16$, $R_\lambda{=}1.0$, $\rho{=}0.6$  & $5^\circ$ & $0/10$ & $32$ & $0/10$ & 17974 \\
$M{=}16$, $R_\lambda{=}1.0$, $\rho{=}0.3$  & $3^\circ$ & $2/10$ & $4$  & $2/10$ & 11109 \\
$\bm{M{=}16}$, $\bm{R_\lambda{=}1.0}$, $\bm{\rho{=}0.3}$ & $\bm{5^\circ}$ &
$\bm{10/10}$ & $\bm{0}$ & $\bm{10/10}$ & $\bm{10214}$ \\
$M{=}16$, $R_\lambda{=}1.0$, LOS only      & $5^\circ$ & $10/10$ & $0$ & $10/10$ & 10299 \\
\bottomrule
\end{tabular}
\end{table}

Table~\ref{tab:rfrun} runs the complete method against the chain. The outcome is sharply
dichotomous: with an adequate aperture and moderate multipath the method clears every emitter,
issues \emph{zero} false absence certificates and certifies every mission at $10214$~s ---
within $1\%$ of the $10299$~s obtained in a pure line-of-sight channel. With
strong multipath no configuration we tried succeeds. The relaxation of
Section~\ref{sec:robust} does not repair it: it absorbs spurious \emph{readings}, whereas the
false absence certificates counted here come from \emph{missed detections} --- the failure mode
that Section~\ref{sec:robust} shows no amount of network redundancy can repair. Two design rules follow directly. The angular bound $\delta$ must be
\emph{calibrated to the channel} --- $\delta=3^\circ$ is too tight and $\delta=5^\circ$ works, in
agreement with the measured $p_{90}$ --- and the antenna aperture must be chosen to resolve the
multipath, because no amount of online reasoning can compensate for a sensor that reports the
wrong angle confidently.


\noindent\textbf{Is the chain an unbiased, efficient implementation?} A simulation can be
self-consistent and still wrong, and with no hardware available we use an external standard
instead: the Cram\'er--Rao lower bound (CRLB). For a single source observed by an $M$-element
uniform circular array of radius $R_\lambda$ wavelengths over $N$ snapshots at per-element SNR
$\gamma$, the deterministic (Slepian--Bangs) bound is
\begin{equation}
\operatorname{var}(\hat\theta)\ \ge\ \frac{1}{N\,\gamma\,M\,(2\pi R_\lambda)^2},
\label{eq:crlb}
\end{equation}
obtained from $\sum_m\sin^2(\theta-\phi_m)=M/2$ and $\sum_m\sin(\theta-\phi_m)=0$. Table~\ref{tab:crlb}
compares \eqref{eq:crlb} with the RMS error of the chain over $4000$ trials.

\begin{table}[t]
\caption{The direction-finding chain against the Cram\'er--Rao lower bound, single source,
line-of-sight. The estimator is efficient to within the $0.1^\circ$ grid on which the MUSIC
spectrum is evaluated.}
\label{tab:crlb}
\centering
\small
\begin{tabular}{lccc}
\toprule
Configuration & CRLB ($^\circ$) & Measured RMS ($^\circ$) & Ratio \\
\midrule
$M{=}8$, $R_\lambda{=}0.35$, $N{=}48$, $12$~dB   & $0.334$ & $0.330$ & $0.99$ \\
$M{=}8$, $R_\lambda{=}0.35$, $N{=}48$, $20$~dB   & $0.133$ & $0.134$ & $1.01$ \\
$M{=}8$, $R_\lambda{=}0.35$, $N{=}256$, $12$~dB  & $0.145$ & $0.149$ & $1.03$ \\
$M{=}8$, $R_\lambda{=}1.00$, $N{=}48$, $12$~dB   & $0.117$ & $0.121$ & $1.04$ \\
$M{=}16$, $R_\lambda{=}1.00$, $N{=}48$, $12$~dB  & $0.083$ & $0.088$ & $1.07$ \\
\bottomrule
\end{tabular}
\end{table}

The measured error tracks the bound to within $7\%$ over a factor of four in aperture and a factor
of eight in SNR, with the small excess explained by the $0.1^\circ$ evaluation grid. The simulated
chain therefore closely tracks the Cram\'er--Rao bound across the tested configurations, which
provides an external consistency check on the bearing model under the assumed line-of-sight
condition.

\noindent\textbf{Parameter provenance.} Finally, none of the physical parameters above is
arbitrary. The multipath-to-line-of-sight amplitude ratio $\rho$ used throughout
Section~\ref{sec:rf} and \ref{sec:ghost} maps to an equivalent Rician factor through $K=20\log_{10}(1/\rho)-10\log_{10}L$, where
$L$ is the number of specular reflections: the model of Section~\ref{sec:rf} uses $L=2$
equal-amplitude reflections, so the multipath-to-line-of-sight \emph{power} ratio is $L\rho^2$.
With $L=2$ the values used correspond to $K=7.4$~dB for $\rho=0.3$, $1.4$~dB for $\rho=0.6$ and
$-1.1$~dB for $\rho=0.8$ --- from a strong line-of-sight channel to one in which the multipath
power exceeds the direct path. Throughout the paper we therefore label the multipath level by $\Gamma_{\mathrm{mp}}=20\log_{10}(1/\rho)$, and the ordinate of Fig.~\ref{fig:reliability} uses $\Gamma_{\mathrm{mp}}$, not $K$. The values used therefore span the range from a strong line-of-sight
(suburban, low-altitude) channel to a weak one with comparable multipath power, rather than being
chosen to make the method look good. The antenna patterns of Table~\ref{tab:antenna} likewise use
beamwidths and front-to-back ratios that are representative of commercially available Yagi and
patch-array hardware rather than idealized limits. Anchoring the model to a \emph{specific},
measured site and antenna remains future work, and is the main reason we regard the present study
as a methodological contribution whose validity domain is stated explicitly.

This experiment is the closest substitute for hardware that we can offer. It exercises the
complete measurement chain --- propagation, array geometry, sampling, covariance estimation and
spectral estimation --- rather than an abstract noise model, and every parameter in it is a
physical quantity. Replacing the simulated snapshots with samples from a software-defined radio
would leave the rest of the pipeline unchanged.

\subsection{Generalization}

Fig.~\ref{fig:gen} sweeps five axes. With the two-layer network every configuration attains a
cleared fraction of $1.000$ and a certification rate of $1.000$; the visible trend is the
amortization of the census cost, the per-emitter cost falling from $835$~s at $k=10$ to $555$~s at
$k=16$. Sensitivity to the reception-radius distribution is below $0.2\%$ because the reasoning is
uniformly conservative, and widening the angular error bound eightfold (from $\delta=0.25^\circ$
to $2^\circ$) costs less than $1.5\%$ --- the practical justification for choosing $\delta$
conservatively, as Lemma~\ref{lem:outer} requires.

The most informative experiment is a controlled comparison of the two networks under an increasing
fraction of directional emitters (Fig.~\ref{fig:dir}). With the seven-point single-layer ring
network that is perfectly adequate for omni-directional emitters, the cleared fraction falls
monotonically from $1.000$ to $0.783$, $0.547$, $0.328$ and finally $0.116$ as the fraction of
directional emitters rises through $0.25$, $0.5$, $0.75$ and $1$; the certification rate drops to
zero as soon as a single directional emitter is present. The mechanism is exactly
Theorem~\ref{thm:detect}: the ring network does not satisfy \eqref{eq:conv} everywhere, so a
directional emitter whose illuminated half-plane faces away from every network point is never
heard, and its absence can never be certified. Replacing \emph{only} the network --- leaving the
entire decision logic untouched --- restores $1.000$ for both metrics at every directional
fraction, at a cost of $2.8\%$ ($616.7$~s to $634.2$~s per emitter). The dense two-layer network is
thus not a conservative over-provision but the enabling condition for the guarantee.

\begin{figure*}[t]
\centering
\includegraphics[width=\textwidth]{en_dir.pdf}
\caption{Controlled comparison of network structures as the fraction of directional emitters
increases, with the decision logic held fixed. The single-layer ring network that suffices for
omni-directional emitters collapses; the two-layer dense network designed to satisfy the condition
of Theorem~\ref{thm:detect} is insensitive.}
\label{fig:dir}
\end{figure*}


\subsection{Obstacles and non-convex workspaces}
\label{sec:obstacles}

The disk geometry of Section~\ref{sec:problem} is convenient but not essential, and it is worth
being explicit about what changes when it is dropped. The detectability condition itself is
geometric and survives verbatim once the line-of-sight relation is made explicit:

\begin{equation}
G\in\conv(V_G),\qquad
V_G=\{P\in\bP:\ \|P-G\|\le\Rmin,\ \mathrm{LOS}(P,G)\},
\label{eq:visconv}
\end{equation}

where $\mathrm{LOS}(P,G)$ holds when the segment $PG$ meets no obstacle. The certificate of
Theorem~\ref{thm:cert} then reads unchanged, with $V_G$ in place of $S_G$. The corollaries,
however, acquire a new failure mode: an obstacle can block a direction completely, in which case
\emph{adding network points on the blocked side is useless}. Points must be placed inside the
shadow instead.

To quantify this we take the 25-point two-layer network that satisfies \eqref{eq:conv} on the disk up to the residual characterised in Section~\ref{sec:network}, and place it in a workspace containing three rectangular obstacles (two walls and one
block, occupying $16.7\%$ of the disk, so $83.3\%$ of it is free space). The condition is
evaluated on the same common $80$~m grid as the comparison of Section~\ref{sec:ncbounds}, so the
two tables are directly comparable.

\begin{table}[t]
\caption{Effect of obstacles on the detectability condition \eqref{eq:visconv}. Every row is
evaluated identically: the fraction of free-space positions on a common $80$~m grid at which the
exact condition holds ($1576$ positions without obstacles, $1301$ with). This is a pointwise
evaluation on a grid, not the conservative cell certificate of Section~\ref{sec:network}; the
control row reads $1.000$ at grid resolution while still carrying the vanishing residual
described in Section~\ref{sec:network}.}
\label{tab:obstacles}
\centering
\small
\begin{tabular}{lccc}
\toprule
Network & Points & Condition holds & First failure \\
\midrule
Two-layer, no obstacles (control)   & $25$ & $1.000$ & --- \\
Two-layer, with obstacles           & $25$ & $0.587$ & $(-1720,-520)$ \\
Greedy shadow-filling, with obstacles & $52$ & $0.977$ & $(-1240,-200)$ \\
\bottomrule
\end{tabular}
\end{table}

The result is severe: once obstacles are present the disk-optimized network fails the condition at
$41.3\%$ of the free-space positions, because a shadow removes exactly the detections that the
enclosure argument relies on. The loss is nonetheless concentrated rather than diffuse --- the
median failing position lies $260$~m from the nearest obstacle, against $522$~m for the positions
that pass --- so a targeted repair recovers most of it. Starting from the failing positions we
greedily add the candidate that makes the largest number of them satisfy \eqref{eq:visconv},
subject to having line-of-sight to them. Twenty-seven added points raise the passing fraction to
$0.977$; the remaining failures lie in the deepest part of the shadows, where the only remedy is to
place a point inside the shadow itself.

\noindent\textbf{Kinematic feasibility.} The analysis so far treats travel as straight-line
displacement between network points, which is invalid once obstacles are present. We therefore
re-plan the census tour on a \emph{visibility graph}: nodes are the network points plus the
obstacle vertices (inflated by a $40$~m platform radius), and an edge joins two nodes that can see
each other. The shortest collision-free path between consecutive census points is then obtained by
Dijkstra's algorithm.

Three findings follow. First, three of the $25$ points of the disk-optimal network ---
$(500,-866)$, $(1500,-866)$ and $(-1000,0)$ in the coordinates of Fig.~\ref{fig:net} --- fall
\emph{inside} obstacles and are physically unreachable; they must be relocated before the plan is
executed. This is a failure mode that no purely geometric certification would catch, and it is
worth stating explicitly: a network that is optimal for an empty workspace may be physically
infeasible in a cluttered one. Second, on the naive straight-line ordering, two of the $22$
segments cross an obstacle, so the straight-line assumption is not merely optimistic but
incorrect. Third, re-routing through the visibility graph increases the census tour from
$15\,760$~m to $18\,255$~m, a penalty of $15.8\%$, i.e. about $500$~s of the mission.

For completeness, an equivalent formulation is to inflate the obstacles and connect consecutive
waypoints by Dubins curves when a minimum turning radius is imposed; the qualitative conclusion is
unchanged, and the $15.8\%$ penalty is a property of the obstacle geometry rather than of the
path-planning algorithm.

Two conclusions are worth recording. First, the plan is portable: the disk-optimized network is a
usable starting point for an obstructed workspace, and repairing it is cheaper than re-optimizing
from scratch, provided that points falling inside obstacles are relocated. Second, the repair is
not free --- it roughly doubles the network, from $25$ to $52$ points, and therefore roughly
doubles the census cost of Section~\ref{sec:experiments}. Whether a region with heavy occlusion is
worth searching with one platform, or should be partitioned among several, is an application
decision that our analysis now makes quantitative.


\subsection{Real antenna patterns versus the ideal half-plane}
\label{sec:antenna}

The half-plane model of \eqref{eq:sensor} assumes unit gain everywhere in a $180^\circ$ sector and
exactly zero gain elsewhere. Real directional antennas --- Yagi, patch arrays, horns --- have a
main lobe of finite beamwidth, sidelobes, a finite back lobe and nulls. Writing $g(\psi)$ for the
normalised voltage pattern at an angle $\psi$ from boresight, and $R_{\mathrm{ref}}$ for the
boresight range at which the received power equals the receiver sensitivity, detection from a
platform at $P$ occurs exactly when
\begin{equation}
\|P-G\|\le R_{\mathrm{ref}}\,g\bigl(\angle(u(d),P-G)\bigr),
\label{eq:pattern}
\end{equation}
and the detectability condition becomes: \emph{for every boresight $d$, at least one network point
lies in the corresponding region} \eqref{eq:pattern}. With $g$ the ideal half-plane indicator this
reduces to $G\in\conv(S_G)$, so Theorem~\ref{thm:detect} is the special case.

\begin{table}[t]
\caption{Detectability with realistic patterns. Left: numerically verified detectability fraction
--- the fraction of the region on which the 25-point network of Section~\ref{sec:network}
satisfies \eqref{eq:pattern} for \emph{every} boresight. Right: hexagonal lattice spacing
required to restore a detectability fraction of $0.999$.}
\label{tab:antenna}
\centering
\small
\begin{tabular}{lccc}
\toprule
Pattern & 25-pt fraction & Required spacing & Points \\
\midrule
Ideal half-plane (paper baseline)        & $1.000$ & $1000$~m & $85$ \\
Cardioid, back lobe $-20$~dB             & $1.000$ & $1000$~m & $85$ \\
Yagi, HPBW $55^\circ$, F/B $16$~dB, SL $-13$~dB & $0.842$ & $800$~m & $97$ \\
Yagi, HPBW $40^\circ$, F/B $20$~dB, SL $-15$~dB & $0.004$ & $800$~m & $97$ \\
Patch array, HPBW $65^\circ$, F/B $25$~dB       & $0.013$ & $600$~m & $127$ \\
\bottomrule
\end{tabular}
\end{table}

The result, Table~\ref{tab:antenna}, is counter-intuitive and important: \emph{the ideal
half-plane is not the worst case}. A narrow main lobe with low sidelobes is far harder to
guarantee, because the boresight can point \emph{between} network points, whereas a half-plane at
least satisfies the condition from any point in a $180^\circ$ arc. With the paper's 25-point
network a realistic patch pattern satisfies it on only $1.3\%$ of the region; a Yagi with $-13$~dB
sidelobes manages $84\%$ precisely because its sidelobes fill the angular gaps. Restoring the
guarantee requires a $1.25$--$1.65\times$ denser network ($600$--$800$~m spacing,
$97$--$127$ points) than the idealised model predicts.

The engineering conclusion is that the idealisation in \eqref{eq:sensor} is \emph{optimistic} and
must be corrected by the measured pattern of the actual antenna before the network is designed.
This is a one-line substitution --- replace the indicator with $g$ --- but it changes the answer by
a factor we now know how to compute.

\subsection{Spatially correlated multipath: the coherent ghost}
\label{sec:ghost}

Section~\ref{sec:robust} models outliers as independent draws. Real multipath is not independent:
reflections come from fixed buildings and terrain, so the spurious bearing \emph{varies smoothly as
the platform moves and always points at the same physical location}. The consequence is
qualitatively different and worse. Independent outliers are mutually inconsistent, so the evidence
set cannot converge on any of them and stays near the truth. A correlated outlier set is
\emph{mutually consistent}: the bearings agree, the evidence set converges confidently on a
\emph{ghost}, and the platform drives to a place where nothing exists.

We model this directly by replacing, with probability $p_g$, the reported bearing by the bearing
towards a fixed virtual reflector. Table~\ref{tab:ghost} compares this with independent outliers.

\begin{table}[t]
\caption{Correlated multipath (a fixed virtual reflector) versus independent outliers, $10$
episodes each, scenario S2. A coherent ghost is far more damaging at the same rate.}
\label{tab:ghost}
\centering
\small
\begin{tabular}{lcccc}
\toprule
Perturbation & Rate & Cleared & \textbf{False abs.} & Certified \\
\midrule
None                        & ---   & $10/10$ & $0$   & $10/10$ \\
Independent outliers        & $0.30$ & $10/10$ & $0$   & $10/10$ \\
\textbf{Coherent ghost}     & $0.30$ & $\bm{0/10}$ & $\bm{44}$ & $0/10$ \\
Coherent ghost              & $0.70$ & $0/10$ & $82$  & $0/10$ \\
Coherent ghost              & $1.00$ & $0/10$ & $100$ & $0/10$ \\
\bottomrule
\end{tabular}
\end{table}

A ghost rate of $0.3$ already destroys the mission, whereas independent outliers at the same rate
are harmless. We then tested three defences, all of which follow from the same idea --- that a
\emph{failed} clearance is physical evidence contradicting the evidence set: the validation gate of
Section~\ref{sec:robust}, adaptive relaxation of $q$ with reconstruction, and rejection of the
bearing cluster that voted for the refuted position. All three reduce the number of false absence
certificates somewhat ($44\to37$ at $p_g=0.3$) and \emph{none} restores the cleared fraction
($0/10$ in every configuration we tried).

We report this as a negative result and draw the honest conclusion: \textbf{the
$q$-relaxed framework is structurally inadequate for spatially correlated multipath}, because its
central assumption --- that outliers disagree with each other and with the truth --- is exactly
what correlated multipath violates. What is required is an architecture that maintains several
source \emph{hypotheses} and resolves them by visiting them, rather than a single evidence set that
is relaxed when contradicted. We regard this as the most important open problem raised by this
work.

\subsection{Intermittent and moving emitters}
\label{sec:burst}

The certificate rests on a determinism assumption: if an emitter is within range and inside its
illuminated half-plane, it \emph{is} heard. Frequency hopping and burst transmission violate this
directly --- a \texttt{no\_signal} reading may mean ``transmitting elsewhere'' or ``not
transmitting now'' rather than ``not present''. We model a duty cycle $p_{\mathrm{on}}$ and report
the consequences in Table~\ref{tab:duty}.

\begin{table}[t]
\caption{Duty cycle: intermittent transmission breaks the absence certificate. $10$ episodes each,
scenario S2.}
\label{tab:duty}
\centering
\small
\begin{tabular}{lccc}
\toprule
Duty cycle $p_{\mathrm{on}}$ & Cleared & \textbf{False abs.} & Certified \\
\midrule
$1.0$ (continuous) & $10/10$ & $0$  & $10/10$ \\
$0.7$              & $5/10$  & $7$  & $5/10$ \\
$0.4$              & $0/10$  & $36$ & $0/10$ \\
\bottomrule
\end{tabular}
\end{table}

A duty cycle of $0.7$ already produces false absence certificates. The certificate is
deterministic only under continuous transmission, and there is no way to repair this within the
present framework: a source that is silent half the time is observationally equivalent to a
source that is absent, and the certificate cannot distinguish them from a single pass.

The only sound remedy is redundancy in time. If the census is repeated $R$ rounds, each round
occupying a different transmission slot, then a present emitter is missed in all $R$ rounds with
probability at most $(1-p_{\mathrm{on}})^R$ (for independent slots), so
\begin{equation}
\Pr[\text{false absence}]\ \le\ (1-p_{\mathrm{on}})^R .
\label{eq:rounds}
\end{equation}
This is a \emph{probabilistic} statement, not a proof, and it must be reported as such: the
guarantee degrades from ``certified absent'' to ``absent with confidence $1-(1-p_{\mathrm{on}})^R$''.
For $p_{\mathrm{on}}=0.4$, $R=6$ gives $4.7\%$ and $R=10$ gives $0.6\%$, at a cost of $R$ times
the census travel. Moving emitters are worse still: without a bound on emitter speed there is no
valid certificate at all, because the region swept at the start of the mission may already be
stale at the end. A speed bound converts this into a tracking problem in which the certificate
must be re-established over a sliding window, which we leave open.


\subsection{Lower bounds for obstructed workspaces}
\label{sec:ncbounds}

Section~\ref{sec:obstacles} showed how to \emph{repair} a network for an obstructed workspace but
left open whether the repair is anywhere near necessary. We close part of that gap here, and are
explicit about the part that remains open.

Throughout, $W\subseteq\bA$ is the free space, $\partial W$ its boundary, $A_{\mathrm{free}}$ its
area, and $V_G=\{P\in\bP:\|P-G\|\le\Rmin,\ \mathrm{LOS}(P,G)\}$ as in \eqref{eq:visconv}.

\begin{theorem}[Occlusion cannot weaken the arc bound]
\label{thm:ncarc}
Let $\varphi_{\mathrm{vis}}=\max_{P\in W}\varphi_{\mathrm{vis}}(P)$, where
$\varphi_{\mathrm{vis}}(P)$ is the half-angle of the longest arc of the outer boundary
$\partial\bA$ that is simultaneously within $\Rmin$ of $P$ and visible from $P$. Then
\begin{equation}
k\ \ge\ \left\lceil\frac{\pi}{\varphi_{\mathrm{vis}}}\right\rceil
\ \ge\ \left\lceil\frac{\pi}{\varphi_{\max}}\right\rceil=6,
\qquad \varphi_{\max}=\arcsin(\Rmin/R).
\label{eq:ncarc}
\end{equation}
\end{theorem}
\begin{bproof}
Every point of $\partial\bA$ is a candidate emitter position and must be detectable. A network
point can only serve those outer-boundary positions that are both within $\Rmin$ and visible;
inserting obstacles can only remove candidates, so $\varphi_{\mathrm{vis}}(P)\le\varphi_{\max}$
for every $P$, with equality in the unobstructed case. Since the served arcs must cover the whole
circle, $k\varphi_{\mathrm{vis}}\ge\pi$.
\end{bproof}

Obstacles therefore \emph{monotonically strengthen} the bound of Theorem~\ref{thm:klb}; they can
never relax it. Evaluating \eqref{eq:ncarc} on the three-obstacle workspace gives the unobstructed
half-angle $\varphi_{\mathrm{vis}}=33.75^\circ$ --- a full visible arc of $67.5^\circ$ ---
i.e. $k\ge\lceil180/33.75\rceil=6$, unchanged --- the obstacles are placed far
enough from the outer boundary not to occlude it. Placing an obstacle against the boundary, or
adding a second outer shell, would raise the bound immediately.

\begin{theorem}[Route-length bound for arbitrary $W$]
\label{thm:nclen}
Any tour that makes every point of $W$ detectable --- in particular any tour visiting a network
satisfying \eqref{eq:visconv} --- has
\begin{equation}
L\ \ge\ \frac{A_{\mathrm{free}}-\pi\Rmin^2}{2\Rmin}.
\label{eq:nclen}
\end{equation}
\end{theorem}
\begin{bproof}
By Corollary~\ref{cor:spacing} the network covers $W$ in the sense that every point of $W$ lies
within $\Rmin$ of some network point, so the tour sweeps $W$ with disks of radius $\Rmin$. A path
of length $L$ sweeps at most $2\Rmin L+\pi\Rmin^2$ of area; equating this with
$A_{\mathrm{free}}$ gives \eqref{eq:nclen}. Obstacles reduce $A_{\mathrm{free}}$ and therefore
\emph{weaken} this particular bound, which is why it is stated in terms of the free area.
\end{bproof}

For the obstructed workspace of Section~\ref{sec:obstacles},
$A_{\mathrm{free}}=8.48\times10^6$~m$^2$ ($83.3\%$ of the disk) and
\eqref{eq:nclen} gives $L\ge2697$~m.

\noindent\textbf{A counting estimate, and why it is not a bound.} It is tempting to add a third
bound by arguing that enclosure requires three points and hence
$k\ge\lceil 3N/s^\ast\rceil$, where $N$ is the number of cells and $s^\ast$ the largest number of
cells a single point can serve. \emph{This is not valid, and we record it only to explain why we do
not use it.} Enclosure does not require three points: if $G$ lies on the segment joining two
network points, or coincides with one, then $G\in\conv(S_G)$ already holds with two points or with
one. The configurations that must be excluded form the zero-measure locus swept out by the
segments joining pairs of network points, so a rigorous version would need an almost-everywhere
argument and an area integral rather than a count of cell centres. Evaluated as stated on an
$80$~m partition of the obstructed workspace ($N=1301$, $s^\ast=474$) the expression suggests
$k\gtrsim9$; we report this as an order-of-magnitude indication only.

\noindent\textbf{What these bounds do and do not settle.} The two rigorous bounds above apply to any workspace, but they are \emph{weak}: against the constructive solution of
Section~\ref{sec:obstacles} ($52$ points, a $30$~km route) they leave gaps of factors $5.8$ and
$11$. The reason is structural rather than technical, and worth stating plainly. Counting arguments of that form: they compare a total demand against the largest
per-point supply, and they are tight only when the binding constraint is \emph{how much} a point
can serve. In this problem the binding constraint is \emph{in which directions} it can serve:
enclosure fails not because too few points are nearby but because the nearby points happen to lie
in a common half-plane. No argument that counts incidences can see this, which is why the
covering-type bounds of Section~\ref{sec:bounds} were tight in the convex case (where a single
arc-length count is essentially the whole story) and are loose here.

Closing the gap requires an \emph{angular} relaxation --- for instance a linear program over
candidate points with one constraint per (cell, direction) pair, whose dual would expose the
angular bottlenecks that our counting bound averages away. We regard this, together with the
correlated-multipath problem of Section~\ref{sec:ghost}, as the two clearly identified open
problems of this work.

One quantitative conclusion does survive, and it is favourable to the repair strategy. On the
common evaluation grid of $1301$ free-space positions --- the same evaluation as
Table~\ref{tab:obstacles}, so the two tables are complementary by construction --- the comparison
is:

\begin{center}
\small
\begin{tabular}{lcc}
\toprule
Network & Points & Positions failing \\
\midrule
Disk-optimal two-layer              & $25$  & $41.3\%$ \\
Disk-optimal $+$ greedy repair      & $52$  & $\bm{2.3\%}$ \\
Uniform hexagonal lattice, $600$~m  & $79$  & $12.45\%$ \\
Uniform hexagonal lattice, $500$~m  & $112$ & $6.07\%$ \\
Uniform hexagonal lattice, $450$~m  & $129$ & $7.61\%$ \\
\bottomrule
\end{tabular}
\end{center}

The greedy repair is thus substantially more efficient \emph{per point} than the best uniform
lattice we could find: $52$ targeted points outperform $112$ uniform ones. Refining a uniform
lattice is moreover not monotone --- $129$ points at $450$~m spacing do \emph{worse} than $112$ at
$500$~m, because a change of phase moves the points relative to the shadows --- which is precisely
the effect a uniform construction cannot control and a greedy one can. This justifies the repair
as a practical answer even though it is not optimal.

The only rigorous point-count result is the lower bound $k\ge6$. The counting estimate suggests a
practical scale nearer $9$, and our $52$-point repair attains $97.7\%$ of the workspace under the
evaluation of Table~\ref{tab:obstacles}. We therefore do \emph{not} claim a rigorous upper bound
on the number of points needed for exact full-workspace certification: the constructive end is a
sufficient design, not a proven minimum. The gap between the two is itself the honest measure of
how much of the non-convex problem remains open, and closing it requires the angular relaxation
described above.


\subsection{Certificate reliability as a function of the sensing chain}
\label{sec:reliability}

The results so far answer the question a designer of the \emph{method} asks: does the strategy
work? This section answers the question a designer of the \emph{sensor} asks, which is the one that
determines whether the certificate can be trusted at all: \emph{given a channel and a direction
finder, what is the probability that the certificate is wrong?} We sweep the two parameters that
characterise the sensing chain --- the aperture $M R_\lambda$ and the per-path multipath level $\Gamma_{\mathrm{mp}}=20\log_{10}(1/\rho)$ (see the parameter-provenance paragraph below for its relation to the Rician $K$-factor) --- together with the
per-element SNR and the transmitter duty cycle, and record two quantities: the \emph{false
certificate rate}, i.e. the fraction of emitters wrongly certified absent, which is the only
irrecoverable error, and the \emph{certified-completion rate}, i.e. the fraction of missions that
terminate with a proof. Each cell is $20$ episodes of scenario S2.

\begin{figure*}[t]
\centering
\includegraphics[width=\textwidth]{reliability_map.pdf}
\caption{Reliability of the certificate as a function of the sensing chain. (a)~False-certificate
rate over the (aperture, multipath-level $\Gamma_{\mathrm{mp}}=20\log_{10}(1/\rho)$) plane; the dashed line is the $1\%$ contour. (b)~The
complementary certified-completion rate. (c)~Dependence on per-element SNR at five apertures;
note the near-flatness. (d)~Dependence on the transmitter duty cycle.}
\label{fig:reliability}
\end{figure*}

\begin{table}[t]
\caption{Selected operating points from Fig.~\ref{fig:reliability}a--b. A certificate worth
trusting requires both a low false-certificate rate and a non-negligible certified-completion rate;
the two conditions together define a frontier, not a point.}
\label{tab:reliability}
\centering
\small
\begin{tabular}{cccc}
\toprule
Aperture $MR_\lambda$ & $\Gamma_{\mathrm{mp}}$ (dB) & False cert. & Certified \\
\midrule
$2.8$  & $3$  & $0.60$ & $0.00$ \\
$4.8$  & $7$  & $0.28$ & $0.00$ \\
$14.4$ & $7$  & $0.07$ & $0.10$ \\
$14.4$ & $9$  & $0.02$ & $0.65$ \\
$19.2$ & $7$  & $0.01$ & $0.90$ \\
$19.2$ & $9$  & $0.00$ & $1.00$ \\
\bottomrule
\end{tabular}
\end{table}

Three conclusions emerge, and the second is the one we consider most useful.

\noindent\textbf{(i) Reliability has a frontier, and it is sharp.} The $1\%$ contour in
Fig.~\ref{fig:reliability}a runs diagonally: a certificate that is reliable in a strong
line-of-sight channel becomes worthless in a channel with comparable multipath power unless the
aperture is increased. At $\Gamma_{\mathrm{mp}}=7$~dB (equivalently $K\approx4$~dB), moving from an aperture of $4.8$ to $19.2$ reduces the false-certificate rate from $0.28$ to $0.01$; at $\Gamma_{\mathrm{mp}}=3$~dB (equivalently $K\approx0$~dB) no aperture we tested is safe.

\noindent\textbf{(ii) SNR is not the control variable.} Fig.~\ref{fig:reliability}c is, to our
knowledge, the most practically useful panel of this paper. Raising the per-element SNR from $6$ to
$20$~dB --- more than a factor of six in power --- leaves the false-certificate rate essentially
unchanged at apertures of $1.6$, $2.8$ and $4.8$, because the bearing error is dominated by
multipath \emph{bias} rather than by noise. Only aperture moves the frontier. The design implication
is direct and, we suspect, counter-intuitive to practitioners accustomed to buying sensitivity:
\emph{integrate longer is not the answer; make the array bigger is.}

\noindent\textbf{(iii) Intermittent transmission is a separate, steeper axis.}
Fig.~\ref{fig:reliability}d shows the duty cycle acting on the certificate independently of the
array: at an aperture of $19.2$ and $\Gamma_{\mathrm{mp}}=7$~dB, where continuous transmission gives a $0.7\%$
false-certificate rate and $90\%$ certified completion, a duty cycle of $0.4$ degrades these to
$25\%$ and $0\%$. No antenna can recover this, because the information is simply absent from the
measurements; only time redundancy can, at the cost of the probabilistic degradation
\eqref{eq:rounds}.

Together, Fig.~\ref{fig:reliability} converts the paper's central claim from a binary statement
into a design procedure: choose the aperture that places the operating point inside the $1\%$
contour for the channel in which the system will actually be deployed, verify the duty cycle
against \eqref{eq:rounds}, and only then trust the certificate. We regard this map --- rather than
any single number in the tables above --- as the most transferable result of the paper.

\subsection{Competitiveness against an omniscient planner}

To put the cost in perspective we compare against an \emph{omniscient reference plan} that knows
all emitter positions in advance and therefore needs no measurements at all: its cost is the open
travelling-salesman tour through the true emitter positions divided by the speed, plus the fixed
neutralization time.

Two remarks on how this quantity may be read. First, it is an \emph{upper} bound on the cost of an
optimal omniscient plan, not a lower bound on the cost of an arbitrary strategy: an omniscient
planner needs only to enter the $r_{\mathrm{clr}}=20$~m neighbourhood of each emitter, so it can
beat the tour through the exact centres. Consequently the ratio we report is a \emph{lower} bound
on the true price of not knowing the emitter positions, and the true gap is at least as large.
Second, the comparison remains informative because the reference plan is a concrete, achievable
omniscient strategy rather than an abstract bound.

The proposed method requires $2.43\pm0.24$ times the reference cost in S1 and $4.68\pm0.62$ in S2. The excess is entirely the price
of discovery --- $58\%$ of the mission in S1 and $78\%$ in S2 --- which is consistent with the
observation that removal itself is cheap ($5$~s per emitter) while acquiring the information
needed to locate the emitters dominates the budget.


\subsection{Behaviour under non-ideal channels}
\label{sec:outlier}

To probe the concern that motivates Section~\ref{sec:robust}, we augment the sensor model with two
impairments of a real radio channel: a \emph{spurious bearing} with probability $p_{\mathrm{out}}$
(a \texttt{no\_signal} outcome is replaced by a \texttt{direction} reading at a random angle, as
multipath and sidelobes would produce) and a \emph{missed detection} with probability
$p_{\mathrm{miss}}$ (a detection is replaced by \texttt{no\_signal}, as shadowing would produce).
Both are generated by deterministic per-position, per-channel streams so that episodes remain
reproducible. Table~\ref{tab:outlier} reports $12$ episodes of scenario S2 for each combination of
$p_{\mathrm{out}}$ and the relaxation parameter $q$, with $p_{\mathrm{miss}}=0$.

\begin{table}[t]
\caption{Non-ideal channels. Each row is $12$ episodes; \emph{false absence} counts emitters that
were wrongly certified absent, which is the only unrecoverable error.}
\label{tab:outlier}
\centering
\small
\begin{tabular}{cccccc}
\toprule
$p_{\mathrm{out}}$ & $q$ & Cleared & \textbf{False absence} & Certified & Time (s) \\
\midrule
$0.00$ & 0 & 1.000 & 0 & 1.000 & 9044 \\
$0.02$ & 0 & 0.000 & \textbf{79} & 0.000 & 9115 \\
$0.05$ & 0 & 0.000 & \textbf{115} & 0.000 & 8381 \\
$0.10$ & 0 & 0.000 & \textbf{158} & 0.000 & 5938 \\
\midrule
$0.00$ & 2 & 0.917 & 0 & 0.917 & 15723 \\
$0.02$ & 2 & 0.167 & \textbf{0} & 0.000 & 23254 \\
$0.05$ & 2 & 0.000 & \textbf{0} & 0.000 & 21501 \\
$0.10$ & 2 & 0.000 & \textbf{0} & 0.000 & 16713 \\
\midrule
$0.00$ & 4 & 0.750 & 0 & 0.750 & 33809 \\
$0.02$ & 4 & 0.833 & \textbf{0} & 0.000 & 44982 \\
$0.05$ & 4 & 0.417 & \textbf{0} & 0.000 & 41434 \\
$0.10$ & 4 & 0.250 & \textbf{0} & 0.000 & 34439 \\
\bottomrule
\end{tabular}
\end{table}

Three results stand out. First, the \emph{hard} outer approximation ($q=0$) is not merely
degraded but unsafe: a spurious-bearing rate of only $2\%$ already produces $79$ false absence
certificates and drives the cleared fraction to zero. This is the dominant failure mode of a certifying
system, and it is invisible in an ideal-channel evaluation.
Second, separating the two absence paths and using the $q$-relaxed intersection removes the error
completely: at $q\ge2$ the number of false absence certificates is zero at every impairment level
we tested. Correctness is restored. Third, and less comfortably, completeness is not: the
certification rate collapses to zero, because with spurious bearings no channel ever satisfies the
zero-detection condition of Theorem~\ref{thm:cert}, and by
Proposition~\ref{prop:r1} the networks we use cannot invoke the redundancy condition of
Theorem~\ref{thm:redundancy}. The mission remains \emph{safe} but no longer \emph{decidable}.

The practical conclusion is a design rule rather than an algorithmic one: if spurious readings
are expected, the network must be built with detection redundancy $r\ge q+1$ \emph{before} the
mission, because no amount of online reasoning can recover a guarantee that the geometry does not
provide. Together with the results of Section~\ref{sec:ablation}, this defines a clear
safety--efficiency frontier: the certificate is what makes the mission terminate early (a factor of
$4.8$), the relaxation is what makes it correct under outliers, and the redundancy $r$ is what
makes it complete again --- at a cost in network density that we have not yet optimized.

\subsection{Robustness}
Three additional stress tests were performed. Changing the quantization scale that defines
``the same location'' for the deterministic error from $0.1$~m to $10$~m changes the per-emitter
cost by at most $0.8\%$ in S1 and $0.1\%$ in S2, confirming that no conclusion depends on this
modelling choice. Imposing a minimum separation of $50$~m between emitters changes the cost by
$0.3\%$. Finally, forcing every reception radius to the worst case of $1000$~m increases the S1
cost from $370.5$~s to $399.5$~s ($+7.8\%$), in line with the conservative use of $\Rmin$
throughout.

%=====================================================================
\section{Discussion}
\label{sec:discussion}

\subsection{Why certification changes the problem}
The classical formulation of such missions optimizes a scalar trade-off between the probability of
detection and the time spent. Our results show that when the emitters have unknown directionality this is not merely a modelling preference but a structural necessity: a non-detection
carries no information about distance, so the ``probability of having missed something'' cannot be
driven to zero by measuring more densely in the same place. What restores decidability is a
property of the \emph{geometry of the network}: an emitter of arbitrary orientation is detectable
exactly when the network surrounds it. This converts an estimation problem into a verification
problem, and the ablation quantifies the payoff at a factor of $4.8$ in mission time in S2
(in S1 the raster path already resolves empty channels, so the certificate changes nothing there).

\subsection{Limitations}
\begin{itemize}
\item \textbf{Simulation only; no hardware validation.} The study is entirely simulation-based.
The sensor and timing models are taken from a documented specification and reproduced faithfully,
and the physical layer is validated against the Cram\'er--Rao bound (Section~\ref{sec:rf}), but no
physical experiment was performed. A bench demonstration with a software-defined radio and a
single direction-finding antenna, or a small wheeled platform in an outdoor field, would test the
modelling assumptions of Section~\ref{sec:problem} directly and is the single most valuable
extension of this work.
\item \textbf{Known error bound.} The evidence-set construction requires the angular error bound
$\delta$ to be known. If the true error exceeds $\delta$ the outer-approximation invariant of
Lemma~\ref{lem:outer} is violated and the certificate becomes unsound. In practice $\delta$ should
be estimated conservatively; our generalization experiment shows that the cost of using a bound
eight times larger than necessary is below $1.5\%$ in S2, which suggests that erring on the safe
side is cheap.
\item \textbf{Conservative treatment of the reception radius.} Reasoning with $\Rmin$ everywhere
is safe but leaves performance on the table when the true radii are large. Estimating the radii
online from repeated detections is an obvious extension.
\item \textbf{Propagation model.} The channel is modelled as a point emitter with a half-plane
radiation pattern, a bounded deterministic bearing error and, in
Section~\ref{sec:outlier}, independent spurious bearings and missed detections. Real channels add
multipath fading, shadowing and non-ideal sidelobes whose statistics are site-specific. Our
robustness analysis covers the two qualitative effects that matter for the certificate --- a false
detection and a lost detection --- but a calibrated channel model, and in particular the joint
statistics of correlated outliers along a trajectory, remains to be studied.
\item \textbf{Obstacles and non-convex workspaces.} The detectability condition and the
certificate extend to any compact workspace once the visibility relation is made explicit,
\eqref{eq:conv} becoming the enclosure condition \eqref{eq:visconv} of
Section~\ref{sec:obstacles}; the network construction of Section~\ref{sec:network} is specific to
the disk, but Section~\ref{sec:obstacles} shows that it can be repaired in place by greedy
shadow-filling at a cost of roughly a doubling of the network. The bounds of
Section~\ref{sec:bounds} --- the arc argument and the swept-area argument --- do not carry over to
a non-convex region and would have to be re-derived; this is the part of the present work that is
most clearly incomplete.
\end{itemize}

\subsection{Future work}
The most immediate open problem is the one identified in
Section~\ref{sec:robust}: constructing a network of \emph{minimum tour length} that attains a
prescribed detection redundancy $r$, so that a mission can be simultaneously safe and decidable
under a bounded number of spurious readings. A related question is whether redundancy can be
bought more cheaply by moving the outer ring further out --- where a single point covers a wider
angular sector for boundary emitters --- than by densifying it in place.
Three further directions appear promising. First, \emph{adaptive network refinement}: rather than
fixing the network in advance, points can be added where the detectability condition is closest to
being violated, which may reduce the number of required measurements below the uniform-lattice
bound. Second, \emph{multi-robot extension}: the network is a set of locations, so the guarantee
carries over to a team, and the tour becomes a vehicle-routing problem; the certificate then
requires that the union of the teams' measurement points satisfies \eqref{eq:conv}. Third,
\emph{moving emitters}: the evidence set becomes a time-indexed set and the update rules acquire a
prediction step, at the cost of losing the static absence certificate.

%=====================================================================
\section{Conclusion}
We studied the problem of finding and neutralizing an unknown number of radio emitters with
unknown directionality using a single bearing-only mobile sensor, with the additional
requirement that the platform be able to certify that none has been missed. We derived a
necessary and sufficient condition for a measurement network to detect an emitter of arbitrary
orientation, namely that the emitter lies in the convex hull of the network points within the
worst-case reception radius, and showed that this implies the network must extend beyond the region
of interest. Guided by the condition we formulated the network design as a covering-salesman
problem, proved a point-count lower bound, and obtained a two-layer construction that satisfies the detectability condition --- analytically almost everywhere and numerically on the residual with a $15^\circ$ margin --- with a route $42\%$ shorter than the natural single-layer alternative. The resulting
certifiable search strategy clears $100\%$ of emitters on 200 Monte-Carlo episodes at
$359.7\pm44.6$~s per emitter; in a separate $30$-seed comparison the strongest of four baselines
clears $51\%$ with $12.9$ times the budget and none
certifies completion in the directional scenario. Ablations attribute the largest single effect to
the completion certificate, which does not change what is cleared but, in scenario S2, reduces
mission time by a factor of $4.8$.

%=====================================================================
\section*{CRediT authorship contribution statement}

\textbf{Changxuan Cao:} Conceptualization, Methodology, Software, Validation, Formal analysis, Investigation, Data curation, Writing -- original draft, Writing -- review and editing, Visualization. \textbf{Zexu Ouyang:} Methodology, Software, Validation, Formal analysis, Writing -- review and editing.

\section*{Funding}

This research did not receive any specific grant from funding agencies in the public, commercial, or not-for-profit sectors.

\section*{Declaration of competing interest}

The authors declare that they have no known competing financial interests or personal relationships that could have appeared to influence the work reported in this paper.

\section*{Declaration of generative AI and AI-assisted technologies in the writing process}

During the preparation of this work the authors used a large language model assistant in order
to check the internal consistency of the manuscript, to verify that the numerical claims match
the generated data files, and to assist with language editing. After using this tool, the
authors reviewed and edited the content as needed and take full responsibility for the content
of the publication.

\section*{Data availability}

The simulator, the search-and-neutralize implementation, the experiment scripts and the generated data that support the findings of this study are openly available at \url{https://github.com/ccx111145/b.t.measure}.

\begin{thebibliography}{99}

\bibitem{torrieri1984}
D.~J. Torrieri, ``Statistical theory of passive location systems,''
\emph{IEEE Trans. Aerosp. Electron. Syst.}, vol.~AES-20, no.~2, pp.~183--198, 1984.

\bibitem{gavish1992}
M.~Gavish and A.~J. Weiss, ``Performance analysis of bearing-only target location algorithms,''
\emph{IEEE Trans. Aerosp. Electron. Syst.}, vol.~28, no.~3, pp.~817--828, 1992.

\bibitem{bishop2010}
A.~N. Bishop, B.~Fidan, B.~D. O. Anderson, K.~Do\u{g}an\c{c}ay, and P.~N. Pathirana,
``Optimality analysis of sensor--target localization geometries,''
\emph{Automatica}, vol.~46, no.~3, pp.~479--492, 2010.

\bibitem{martinez2006}
S.~Mart\'inez and F.~Bullo, ``Optimal sensor placement and motion coordination for target
tracking,'' \emph{Automatica}, vol.~42, no.~4, pp.~661--668, 2006.

\bibitem{chan1994}
Y.~T. Chan and K.~C. Ho, ``A simple and efficient estimator for hyperbolic location,''
\emph{IEEE Trans. Signal Process.}, vol.~42, no.~8, pp.~1905--1915, 1994.

\bibitem{ho2004}
K.~C. Ho and W.~Xu, ``An accurate algebraic solution for moving source location using TDOA and
FDOA measurements,'' \emph{IEEE Trans. Signal Process.}, vol.~52, no.~9, pp.~2453--2463, 2004.

\bibitem{choset2001}
H.~Choset, ``Coverage for robotics --- a survey of recent results,''
\emph{Annals of Mathematics and Artificial Intelligence}, vol.~31, pp.~113--126, 2001.

\bibitem{galceran2013}
E.~Galceran and M.~Carreras, ``A survey on coverage path planning for robotics,''
\emph{Robotics and Autonomous Systems}, vol.~61, no.~12, pp.~1258--1276, 2013.

\bibitem{choset1998}
H.~Choset and P.~Pignon, ``Coverage path planning: The boustrophedon cellular decomposition,''
in \emph{Field and Service Robotics}, 1998, pp.~203--209.

\bibitem{golden2012}
B.~Golden, Z.~Naji-Azimi, S.~Raghavan, M.~Salari, and P.~Toth, ``The generalized covering
salesman problem,'' \emph{INFORMS J. Comput.}, vol.~24, no.~4, pp.~534--553, 2012.

\bibitem{bezdek1983}
K.~Bezdek, ``\"Uber einige Kreis\"uberdeckungen,'' \emph{Beitr\"age zur Algebra und Geometrie},
vol.~14, pp.~7--14, 1983.

\bibitem{nica2011}
A.~Nica, ``Covering a disk with smaller disks,'' M.Math. thesis, University of Waterloo, 2011.

\bibitem{jung1901}
H.~Jung, ``\"Uber die kleinste Kugel, die eine r\"aumliche Figur einschliesst,''
\emph{J. Reine Angew. Math.}, vol.~123, pp.~241--257, 1901.

\bibitem{welzl1991}
E.~Welzl, ``Smallest enclosing disks (balls and ellipsoids),'' in \emph{New Results and New
Trends in Computer Science}, 1991, pp.~359--370.

\bibitem{vergassola2007}
M.~Vergassola, E.~Villermaux, and E.~D. Chialvo, ```Infotaxis' as a strategy for searching
without gradients,'' \emph{Nature}, vol.~445, pp.~406--409, 2007.

\bibitem{stone2007}
L.~D. Stone, \emph{Theory of Optimal Search}, 2nd ed. INFORMS, 2007.

\bibitem{chung2011}
T.~H. Chung, G.~A. Hollinger, and V.~Isler, ``Search and pursuit--evasion in mobile robotics: A
survey,'' \emph{Autonomous Robots}, vol.~31, no.~4, pp.~299--316, 2011.

\bibitem{azuma2012}
S.~Azuma, M.~S. Sakar, and G.~J. Pappas, ``Stochastic source seeking by mobile robots,''
\emph{IEEE Trans. Autom. Control}, vol.~57, no.~9, pp.~2308--2321, 2012.

\bibitem{guvensan2011}
M.~A. Guvensan and A.~G. Yilmaz, ``On coverage issues in directional sensor networks: A survey,''
\emph{Ad Hoc Networks}, vol.~9, no.~7, pp.~1238--1255, 2011.

\bibitem{milanese1991}
M.~Milanese and A.~Vicino, ``Optimal estimation theory for dynamic systems with set membership
uncertainty: An overview,'' \emph{Automatica}, vol.~27, no.~6, pp.~997--1009, 1991.

\bibitem{jaulin2001}
L.~Jaulin, M.~Kieffer, O.~Didrit, and \'E.~Walter, \emph{Applied Interval Analysis}.
London, U.K.: Springer, 2001.

\bibitem{chisci1996}
L.~Chisci, A.~Garulli, and G.~Zappa, ``Recursive state bounding by parallelotopes,''
\emph{Automatica}, vol.~32, no.~7, pp.~1049--1055, 1996.

\bibitem{alamo2005}
T.~Alamo, J.~M. Bravo, and E.~F. Camacho, ``Guaranteed state estimation by zonotopes,''
\emph{Automatica}, vol.~41, no.~6, pp.~1035--1043, 2005.

\bibitem{combastel2003}
C.~Combastel, ``A state bounding observer based on zonotopes,'' in \emph{Proc. European Control
Conf. (ECC)}, 2003.

\bibitem{scott2016}
J.~K. Scott, D.~M. Raimondo, G.~R. Marseglia, and R.~D. Braatz, ``Constrained zonotopes: A new
tool for set-based estimation and fault detection,'' \emph{Automatica}, vol.~69, pp.~405--413,
2016.

\bibitem{tan2021}
C.~S. Tan, R.~Mohd-Mokhtar, and M.~R. Arshad, ``A comprehensive review of coverage path planning
in robotics using classical and heuristic algorithms,'' \emph{IEEE Access}, vol.~9,
pp.~119\,270--119\,305, 2021.

\bibitem{cabreira2019}
T.~M. Cabreira, L.~B. Brisolara, and P.~R. Ferreira~Jr., ``Survey on coverage path planning with
unmanned aerial vehicles,'' \emph{Drones}, vol.~3, no.~1, art.~4, 2019.

\bibitem{choset2005}
H.~Choset, K.~M. Lynch, S.~Hutchinson, G.~Kantor, W.~Burgard, L.~E. Kavraki, and S.~Thrun,
\emph{Principles of Robot Motion: Theory, Algorithms, and Implementations}. Cambridge, MA, USA:
MIT Press, 2005.

\bibitem{robin2016}
C.~Robin and S.~Lacroix, ``Multi-robot target detection and tracking: Taxonomy and survey,''
\emph{Autonomous Robots}, vol.~40, no.~4, pp.~729--760, 2016.

\bibitem{khan2018}
A.~Khan, B.~Rinner, and A.~Cavallaro, ``Cooperative robots to observe moving targets: Review,''
\emph{IEEE Trans. Cybern.}, vol.~48, no.~1, pp.~187--198, 2018.

\bibitem{grocholsky2006}
B.~Grocholsky, J.~Keller, V.~Kumar, and G.~Pappas, ``Cooperative air and ground surveillance,''
\emph{IEEE Robot. Autom. Mag.}, vol.~13, no.~3, pp.~16--25, 2006.

\bibitem{gonzalez1985}
T.~F. Gonzalez, ``Clustering to minimize the maximum intercluster distance,''
\emph{Theoretical Computer Science}, vol.~38, pp.~293--306, 1985.

\bibitem{beardwood1959}
J.~Beardwood, J.~H. Halton, and J.~M. Hammersley, ``The shortest path through many points,''
\emph{Proc. Cambridge Philos. Soc.}, vol.~55, no.~4, pp.~299--327, 1959.

\bibitem{applegate2006}
D.~L. Applegate, R.~E. Bixby, V.~Chv\'atal, and W.~J. Cook, \emph{The Traveling Salesman
Problem: A Computational Study}. Princeton, NJ, USA: Princeton Univ. Press, 2006.

\bibitem{koopman1980}
B.~O. Koopman, \emph{Search and Screening: General Principles with Historical Applications}.
New York, NY, USA: Pergamon Press, 1980.

\bibitem{thrun2005}
S.~Thrun, W.~Burgard, and D.~Fox, \emph{Probabilistic Robotics}. Cambridge, MA, USA: MIT Press,
2005.

\bibitem{schmidt1986}
R.~O. Schmidt, ``Multiple emitter location and signal parameter estimation,''
\emph{IEEE Trans. Antennas Propag.}, vol.~34, no.~3, pp.~276--280, 1986.

\bibitem{krim1996}
H.~Krim and M.~Viberg, ``Two decades of array signal processing research: The parametric
approach,'' \emph{IEEE Signal Process. Mag.}, vol.~13, no.~4, pp.~67--94, 1996.

\bibitem{preparata1985}
F.~P. Preparata and M.~I. Shamos, \emph{Computational Geometry: An Introduction}. Springer, 1985.

\bibitem{megiddo1983}
N.~Megiddo, ``Linear-time algorithms for linear programming in $R^3$ and related problems,''
\emph{SIAM J. Comput.}, vol.~12, no.~4, pp.~759--776, 1983.

\bibitem{boyd2004}
S.~Boyd and L.~Vandenberghe, \emph{Convex Optimization}. Cambridge Univ. Press, 2004.

\bibitem{toth1998}
G.~T\'oth, ``Covering a plane convex body by smaller homothetic copies,''
\emph{Discrete Mathematics}, vol.~187, no.~1--3, pp.~281--285, 1998.

\bibitem{popovic2024}
M.~Popovi\'c, J.~Ott, J.~R\"uckin, and M.~J.~Kochenderfer, ``Learning-based methods for adaptive
informative path planning,'' \emph{Robotics and Autonomous Systems}, vol.~179, art.~104727, 2024,
doi:10.1016/j.robot.2024.104727.

\bibitem{kudyba2024}
P.~S. Kudyba, Q.~Lu, and H.~Sun, ``Bayesian optimization for fast radio mapping and
localization with an autonomous aerial drone,'' in \emph{Proc. IEEE Veh. Technol. Conf.
(VTC2024-Fall)}, 2024, doi:10.1109/VTC2024-Fall63153.2024.10757589.

\bibitem{vashisth2026}
A.~Vashisth, M.~Kulshrestha, D.~Conover, and A.~Bera, ``Scalable multi-robot informative path
planning for target mapping via deep reinforcement learning,'' \emph{IEEE Robot. Autom. Lett.},
vol.~11, no.~3, pp.~3414--3421, 2026.

\bibitem{ibenthal2023}
J.~Ibenthal, L.~Meyer, H.~Piet-Lahanier, and M.~Kieffer, ``Localization of partially hidden
moving targets using a fleet of UAVs via bounded-error estimation,'' \emph{IEEE Trans.
Robotics}, vol.~39, no.~6, pp.~4211--4229, 2023.

\bibitem{peng2024}
J.~Peng, Q.~Wang, B.~Jin, Y.~Zhang, and K.~Lu, ``Trajectory optimization to enhance
observability for bearing-only target localization and sensor bias calibration,''
\emph{Biomimetics}, vol.~9, no.~9, art.~510, 2024, doi:10.3390/biomimetics9090510.

\bibitem{kim2025}
S.~Kim, S.~Moon, I.~Petrunin, H.-S.~Shin, and S.~Khattak, ``Autonomous robotic radio source
localization via a novel Gaussian mixture filtering approach,'' in \emph{Proc. Int. Conf.
Inf. Fusion (FUSION)}, 2025, doi:10.23919/FUSION65864.2025.11124026.

\end{thebibliography}

\end{document}

